\documentclass[11pt, reqno]{amsart}
\pagestyle{plain}
\raggedbottom
\usepackage[dvipsnames]{xcolor}
\usepackage[utf8]{inputenc}
\usepackage{amssymb}
\usepackage{amsmath}
\usepackage{amsthm}
\usepackage{amsfonts}
\usepackage[alphabetic]{amsrefs} % for citing inside environments, not compatible with ``cite'' 
\usepackage{mathtools}
\usepackage{fullpage}
\usepackage{float}
\usepackage{latexsym}
\usepackage{bbm}
\usepackage{microtype}
\usepackage{fancyvrb}
\usepackage{mathtools}
\usepackage{tikz, tikz-cd}
\usepackage{enumerate}
\usetikzlibrary{decorations.pathreplacing}
\usepackage{hyperref}
\hypersetup{colorlinks=true, linkcolor=blue, citecolor=magenta, filecolor=magenta, urlcolor=magenta}
\usepackage{bm}
\usepackage{comment}
\allowdisplaybreaks[1]


\usepackage[colorinlistoftodos]{todonotes}
\newcommand{\John}[1]{\todo[size=\tiny,inline,color=blue!30]{#1 \\ \hfill --- John}}
\newcommand{\Nicolle}[1]{\todo[size=\tiny,inline,color=red!30]{#1 \\ \hfill --- Nicolle}}
\newcommand{\Pablo}[1]{\todo[size=\tiny,inline,color=green!30]{#1 \\ \hfill --- Pablo}}

\newtheorem*{theorem*}{Theorem}

\newtheorem{theorem}{Theorem}[section]
\newtheorem{lemma}[theorem]{Lemma}
\newtheorem{proposition}[theorem]{Proposition}
\newtheorem{corollary}[theorem]{Corollary}
\newtheorem*{corollary*}{Corollary}
\newtheorem{conjecture}[theorem]{Conjecture}
\newtheorem*{conjecture*}{Conjecture}
\newtheorem{question}[theorem]{Question}

\newtheorem*{zconjecture*}{Zabrocki's Conjecture}

\newtheorem{introthm}{Theorem}
\renewcommand{\theintrothm}{\Alph{introthm}}
\newtheorem{introcor}[introthm]{Corollary}
\renewcommand{\theintrocor}{\Alph{introthm}}

\theoremstyle{definition}
\newtheorem{definition}[theorem]{Definition}
\newtheorem{example}[theorem]{Example}
\newtheorem{construction}[theorem]{Construction}

\theoremstyle{remark}
\newtheorem{remark}[theorem]{Remark}

\numberwithin{equation}{section}

\makeatletter
\newtheorem*{rep@theorem}{\rep@title}\newcommand{\newreptheorem}[2]{%
\newenvironment{rep#1}[1]{%
\def\rep@title{\bf #2 \ref{##1}}%
\begin{rep@theorem}}%
{\end{rep@theorem}}}
\makeatother
\newreptheorem{theorem}{Theorem}
\newreptheorem{corollary}{Corollary}

\renewcommand{\emptyset}{\varnothing}
\renewcommand{\sl}{\mathfrak{sl}}
\newcommand{\0}{\emptyset}
\newcommand{\Q}{\mathbb{Q}}
\newcommand{\R}{\mathbb{R}}
\newcommand{\Z}{\mathbb{Z}}
\newcommand{\C}{\mathbb{C}}
\newcommand{\N}{\mathbb{N}}
\renewcommand{\S}{\mathfrak{S}}
\newcommand{\cC}{\mathcal{C}}
\newcommand{\cY}{\mathcal{Y}}
\newcommand{\m}{\underline{m}}
\renewcommand{\i}{\underline{i}}
\renewcommand{\r}{\underline{r}}
\newcommand{\s}{\underline{s}}
\newcommand{\e}{\underline{e}}

\DeclareMathOperator{\Hilb}{Hilb}
\DeclareMathOperator{\Frob}{Frob}
\DeclareMathOperator{\Char}{Char}
\DeclareMathOperator{\Sym}{Sym}
\newcommand{\Rho}{\mathrm{P}}
\DeclareMathOperator{\sgn}{sgn}
\DeclareMathOperator{\triv}{triv}
\DeclareMathOperator{\sort}{sort}
\DeclareMathOperator{\area}{area}
\DeclareMathOperator{\leg}{leg}
\DeclareMathOperator{\bounce}{bounce}
\DeclareMathOperator{\dinv}{dinv}
\DeclareMathOperator{\diag}{diag}
\DeclareMathOperator{\coarm}{\mathsf{coarm}}
\DeclareMathOperator{\coleg}{\mathsf{coleg}}
\DeclareMathOperator{\Hom}{Hom}
\DeclareMathOperator{\Ind}{Ind}
\DeclareMathOperator{\Res}{Res}
\DeclareMathOperator{\id}{id}
\DeclareMathOperator{\Span}{span}
\DeclareMathOperator{\shape}{shape}
\DeclareMathOperator{\SYT}{SYT}
\DeclareMathOperator{\SYBT}{SYBT}
\DeclareMathOperator{\maj}{maj}
\DeclareMathOperator{\fmaj}{fmaj}
\DeclareMathOperator{\soc}{soc}

\DeclareRobustCommand{\qbinom}{\genfrac[]{0pt}{}}

\newcommand{\newword}[1]{\emph{\textbf{#1}}}

%%%%%%%%%%%%%%%%%%%%%%%%%%%%%%%%%%%%%%%%%%%%%%%%%%%%%%%%%%%%%%%%%%%%%%
%%%%%%%%%% TITLE %%%%%%%%%%%%%%%%%%%%%%%%%%%%%%%%%%%%%%%%%%%%%%%%%%%%%
%%%%%%%%%%%%%%%%%%%%%%%%%%%%%%%%%%%%%%%%%%%%%%%%%%%%%%%%%%%%%%%%%%%%%%
%%%%%%%%%%%%%%%%%%%%%%%%%%%%%%%%%%%%%%%%%%%%%%%%%%%%%%%%%%%%%%%%%%%%%%

%%%%%%%%%%%%%%%%%%%%%%%%%%%%%%%%%%%%%%%%%%%%%%%%%%%%%%%%%%%%%%%%%%%%%%
%%%%%%%%%% AUTHORS %%%%%%%%%%%%%%%%%%%%%%%%%%%%%%%%%%%%%%%%%%%%%%%%%%%
%%%%%%%%%%%%%%%%%%%%%%%%%%%%%%%%%%%%%%%%%%%%%%%%%%%%%%%%%%%%%%%%%%%%%%
\author[N. González]{Nicolle González}
\address[González]{Department of Mathematics, 
         University of British Columbia, 
         Vancouver, BC, Canada}
\email{nicolle@math.ubc.ca}

\author[J. Lentfer]{John Lentfer}
\address[Lentfer]{Department of Mathematics, 
         University of California, 
         San Diego, CA, USA}
\email{jlentfer@ucsd.edu}

\author[P.S. Ocal]{Pablo S. Ocal}
\address[Ocal]{Department of Mathematics, 
        University of British Columbia, 
        Vancouver, BC, Canada \newline and 
        Department of Mathematics, 
        Osaka Metropolitan University, 
        Osaka, Japan}
\email{socal@math.ubc.ca, y26177q@omu.ac.jp}

%%%%%%%%%%%%%%%%%%%%%%%%%%%%%%%%%%%%%%%%%%%%%%%%%%%%%%%%%%%%%%%%%%%%%%
\setcounter{tocdepth}{1}

\begin{document}
\title{
Sign components for general bosonic-fermionic coinvariant rings
}

%%%%%%%%%%%%%%%%%%%%%%%%%%%%%%%%%%%%%%%%%%%%%%%%%%%%%%%%%%%%%%%%%%%%%%
%%%%%%%%%% ABSTRACT %%%%%%%%%%%%%%%%%%%%%%%%%%%%%%%%%%%%%%%%%%%%%%%%%%
%%%%%%%%%%%%%%%%%%%%%%%%%%%%%%%%%%%%%%%%%%%%%%%%%%%%%%%%%%%%%%%%%%%%%%
\begin{abstract}
    We prove that the sign-isotypic components of the symmetric group coinvariant rings 
    $R_n^{(k,\ell)} \otimes R_n^{(a,b)}$ and $R_n^{(k+a,\ell+b)}$ 
    are isomorphic if and only if $n \leq 2$ or at least one of $(k,\ell)$ or $(a,b)$ is $(0,0)$ or $(0,1)$.
    Moreover, for $k \geq 1$ the $\nu$-isotypic components $(R_n^{(k,\ell)} \otimes R_n^{(0,1)})^\nu$ and $(R_n^{(k,\ell+1)})^\nu$ 
    are isomorphic if and only if $\nu$ corresponds to the sign representation. 
    We also show that the socle of $R_n^{(k,\ell)}$ always contains the sign component, 
    extending one direction of a result of Haiman for $R_n^{(2,0)}$.
    We settle the \emph{hook components} conjecture of Bergeron (2024) 
    for coinvariant rings,
    and a conjecture of Bergeron (2020) that the multiplicity of the sign character 
    of $R_n^{(1,2)}$ is $3^{n-1}$. 
    Our results also recover several existing results in the literature. 
    Lastly, we extend our results to the hyperoctahedral group ${\mathfrak{B}_n}$ 
    and prove that the determinant-isotypic components of 
    $R_{\mathfrak{B}_n}^{(k,\ell)} \otimes R_{\mathfrak{B}_n}^{(0,1)}$ and $R_{\mathfrak{B}_n}^{(k,\ell+1)}$ 
    are isomorphic. 
\end{abstract}
%%%%%%%%%%%%%%%%%%%%%%%%%%%%%%%%%%%%%%%%%%%%%%%%%%%%%%%%%%%%%%%%%%%%%%

\maketitle
\tableofcontents

%%%%%%%%%%%%%%%%%%%%%%%%%%%%%%%%%%%%%%%%%%%%%%%%%%%%%%%%%%%%%%%%%%%%%%
\section{Introduction}
%%%%%%%%%%%%%%%%%%%%%%%%%%%%%%%%%%%%%%%%%%%%%%%%%%%%%%%%%%%%%%%%%%%%%%

For a finite group $G$ with a graded module $M$, the $G$-coinvariant space of $M$ is
defined as the quotient of $M$ by the ideal of homogeneous 
$G$-invariants supported in positive degree. 
Such quotient spaces can be seen as refined versions of the original modules, 
where redundant copies of irreducible components have been eliminated. 
The \newword{$(k,\ell)$ bosonic-fermionic coinvariant ring} $R_n^{(k,\ell)}$ 
is the symmetric group coinvariant space of the polynomial ring 
$\C_n^{(k,\ell)} \coloneqq \C[\bm{x}^{(1)}, \ldots, \bm{x}^{(k)}] \otimes 
\bigwedge[\bm{\theta}^{(1)}, \ldots, \bm{\theta}^{(\ell)}]$ 
in $k$ commuting (bosonic) and $\ell$ anticommuting (fermionic) alphabets in $n$ letters, 
under the diagonal $\S_n$-action which acts on each alphabet 
by permuting the letters simultaneously. The main goal of this paper is to study 
the sign-isotypic component of $R_n^{(k,\ell)}$. 

The study of bosonic-fermionic coinvariant rings goes back to the work of 
Artin in 1944 \cite{Artin} where the classical coinvariant ring $R_n^{(1,0)}$ 
was shown to have dimension $n!$. 
The work of Leray, Borel, and Chevalley \cites{Leray, Borel, Chevalley} 
proves that $R_n^{(1,0)}$ is isomorphic to the regular representation of $\S_n$ 
and to the cohomology ring of the flag variety. 
In the 1990's, Garsia and Haiman introduced and studied the diagonal coinvariant ring 
$R_n^{(2,0)}$ in a series of landmark papers \cites{GarsiaHaiman1993, Garsia-Haiman, Haiman1994} 
that transformed contemporary algebraic combinatorics with striking conjectures 
concerning the dimension and Frobenius character of this ring. 
The $n!$ conjecture and the $(n+1)^{n-1}$ conjecture, named after 
the conjectural dimension of the Garsia--Haiman module and the conjectural dimension of 
$R_n^{(2,0)}$, were resolved by Haiman almost ten years later by appealing, unexpectedly, 
to the geometry of the Hilbert scheme of points on $\C^2$ \cites{Haiman2002, Haiman-polygraph}. 
Haiman's character formula for $R_n^{(2,0)}$, given by the action of the Macdonald eigenoperator 
$\nabla$ \cite{BergeronGarsia} applied to the symmetric functions $s_{(1^n)} = e_n$, 
precipitated a flurry of work seeking a combinatorial formula for this expression, 
culminating in the \emph{shuffle conjecture} \cite{HHLRU2005}. 
In 2015, this conjecture was resolved by Carlsson and Mellit \cite{CarlssonMellit2018} 
(see also \cites{Mellit, BHMPS-Paths, GSV-shuffle}), introducing in the process the novel 
$\mathbb{A}_{q,t}$ algebra, with far reaching algebraic and geometric 
consequences \cites{CGM-aqt, GGS-calibrated, GGS-elliptic, GGS-smooth}.  
    
Motivated by the rich theory of $R_n^{(2,0)}$, 
analogous questions have been asked for other $(k,\ell)$-bosonic-fermionic coinvariant rings. 
For the \emph{superspace coinvariant ring} $R_n^{(1,1)}$, conjectures of 
N. Bergeron, Colmenarejo, Li, Machacek, Sulzgruber, and Zabrocki around 2018 
and of Sagan and Swanson \cite{SaganSwanson2024}, have since been settled: 
its sign character \cite{SwansonWallach1}, 
its dimension and Hilbert series \cite{RhoadesWilson2023}, 
a monomial basis \cites{Angarone2024}, 
and its Frobenius series \cite{MuraiRhoadesWilson} are now established. 
The definition of the bosonic-fermionic coinvariant rings $R_G^{(k,\ell)}$ extends to any finite complex reflection group $G$ \cite{Bergeron2013}, or more generally any finite group $G \subseteq GL(n)$ \cite{Lentfer-Supersymmetry}, although far less is known at these levels of generality. 

Building on the work of Kim and Rhoades \cite{KimRhoades2022}, Jiang and the second author \cite{JiangLentfer} recently established the bigraded Frobenius series of $R_{\mathfrak{B}_n}^{(0,2)}$ explicitly for the hyperoctahedral group $\mathfrak{B}_n = \Z_2 \wr \S_n$. Bhattacharya, Corteel, Jiang, the second author, and Rhoades \cite{BCJLR} recently established the bigraded Frobenius series of $R_{\Z_r \wr \S_n}^{(1,1)}$ for the wreath product groups $\Z_r \wr \S_n$.

Despite the recent advances, computing the full Frobenius series beyond small values of $k$ and $\ell$ has been extremely difficult, with no known unified approach (the Frobenius series is unknown when $k + \ell \geq 3$). 
A more amenable starting point is to focus on specific isotypic components. 
Sign-isotypic components are of particular interest; they frequently contain rich algebraic, combinatorial, and even topological information. 
The sign-isotypic component of $R_n^{(2,0)}$, for instance, has bigraded Hilbert series equal to the \emph{$(q,t)$-Catalan polynomial} and thus computes the degree zero slice of the Khovanov--Rozansky homology of the $(n,n+1)$ torus knot \cite{Hog17}. 

Recently, the first two authors and Mularczyk \cite{GLM} computed the Hilbert series of the sign-isotypic component of $R_n^{(2,1)}$ by proving that the multiplication map $\pi: R_n^{(2,0)} \otimes R_n^{(0,1)} \to R_n^{(2,1)}$ is an isomorphism on sign-isotypic components, establishing the sign component case of Zabrocki's conjecture \cite{Zabrocki2019} that $R_n^{(2,1)}$ is a module for the Delta theorem \cites{HaglundRemmelWilson2018, DAdderioMellit, BHMPS-Delta}. Combined with Haiman's celebrated result for the character of $R_n^{(2,0)}$, the first two authors and Mularczyk determined that the trigraded Hilbert series of $(R_n^{(2,1)})^{\sgn}$ equals a multiple of the \emph{Schr\"oder polynomial} $S_n(q,t,a)$ and consequently coincides with the full Khovanov--Rozansky homology of the $(n,n+1)$ torus knot \cites{GorskyMellit, Hog17}. 

In this article we extend the main result of \cite{GLM} to all bosonic-fermionic coinvariant rings 
$R_G^{(k,\ell)}$ for $G$ the symmetric or hyperoctahedral groups, and prove that the multiplication map 
$\pi: R_G^{(k,\ell)} \otimes R_G^{(0,1)} \to R_G^{(k,\ell+1)}$ is an isomorphism on the sign or determinant isotypic components, respectively. 
In fact, for $G=\S_n$ and $k \geq 1$ we show that the only isotypic component on which $\pi$ is an isomorphism is sign.
Moreover, when $n\geq 3$ we show that the surjection $R_n^{(k,\ell)} \otimes R_n^{(a,b)} \to R_n^{(k+a,\ell+b)}$ restricts to an isomorphism on sign components for nontrivial values of $(k,\ell), (a,b)$ if and only if one of them equals $(0,1)$. 
As a consequence of our construction, 
we also prove that the socle of $R_n^{(k,\ell)}$ contains its sign-isotypic component, 
generalizing the inclusion $(R_n^{(2,0)})^{\sgn} \subseteq \soc(R_n^{(2,0)})$, 
one half of Haiman's result \cite{Haiman2002}*{Sec.~4} that 
the socle of $R_n^{(2,0)}$ is its sign component. 

Although this work extends the main result in \cite{GLM}, the methods employed herein are different 
and in fact dual in a precise way. The results for $R_n^{(2,1)}$ in \cite{GLM} are obtained 
by passing to the harmonic space and realizing $R_n^{(2,1)}$ as 
the space of diagonal super harmonics $H_n^{(2,1)}$. 
Here, we carry out all computations and results within the coinvariant ring $R_n^{(k,\ell)}$. 
Since $R_n^{(k,\ell)}$ and its harmonic analogue $H_n^{(k,\ell)}$ are dual with respect to 
Rota's apolar form \cite{Haiman1994}*{Sec~1.3} (see also \cite{JiangLentfer}*{App. A} 
and \cite{SwansonWallach1}*{Sec. 2.4}), and hence isomorphic as graded $\S_n$-modules, 
one could dualize the methods of this article to obtain an alternative approach in the harmonic space, 
more closely following \cite{GLM}. 

\subsection{Main results}
As above, let $\C_n^{(k,\ell)}$ consist of $k$ commuting and $\ell$ anticommuting families of $n$ 
variables with $I_n^{(k,\ell)} = \left\langle \left(\C_n^{(k,\ell)}\right)^{\S_n}_+\right\rangle$ 
the ideal of $\S_n$-invariants without constant term. 
It is a classic result of Weyl \cite{Weyl}, generalized in \cites{OrellanaZabrocki, JiangLentfer}, 
that $I_n^{(k,\ell)}$ is generated by the \emph{polarized power sums} $p_{\r;\s}$, see Proposition \ref{prop:ideal-generators}.
The ring $R_n^{(k,\ell)} = \C_n^{(k,\ell)} / I_n^{(k,\ell)}$ is then a $(k+\ell)$-graded 
$\S_n$-module with an $\S_n$-module homomorphism $\pi: R_n^{(k,\ell)} \otimes R_n^{(0,1)} \to R_n^{(k,\ell+1)}$ 
sending $f \otimes g \mapsto fg$ whose kernel is the ideal generated by the polarized power sums 
$p_{\r;\s;1}$ of $\C_n^{(k,\ell+1)}$ involving the additional alphabet $\bm{\theta}^{(\ell+1)}$. 
Thus, to prove $\pi$ is injective on sign components we show the restricted kernel 
$\ker(\pi)^{\sgn} = \sum_{(\r,\s) \neq (\underline{0},\underline{0})} \left(R_n^{(k,\ell)} 
\otimes R_n^{(0,1)}\right)^{\sgn} \cdot p_{\r;\s;1}$ is trivial, see Corollary \ref{cor:kernel-isotypic}.

The results in this paper hinge on the fact that $R_n^{(0,1)}$ is isomorphic to the exterior algebra 
$\bigwedge \C^{n-1}$ and thus decomposes as a graded sum of Specht modules indexed by hook partitions \cite{HaglundSergel}. 
Consequently, computing the sign-isotypic component of any module 
$M \otimes \bigwedge \C^{n-1}$ reduces to understanding the hook components of $M$. 
The tensor product $R_n^{(k,\ell)} \otimes \bigwedge \C^{n-1}$ is graded by homogeneous degree in $\bigwedge \C^{n-1}$. 
Working one exterior degree $d$ at a time, 
this reduction is realized by an explicit isomorphism sending 
$f \in \left(R_n^{(k,\ell)} \otimes \bigwedge\nolimits^d \C^{n-1}\right)^{\sgn}$ to 
$\Psi_f \in \Hom_{\S_n}\left(\bigwedge\nolimits^{n-d-1}\C^{n-1}, R_n^{(k,\ell)}\right)$, see Proposition~\ref{prop:Hom-iso}:
\begin{equation}\label{eq:INTRO-hook-hom-iso}
\begin{aligned}
    \left(R_n^{(k,\ell)} \otimes \bigwedge\nolimits^d \C^{n-1}\right)^{\sgn} 
    &\cong 
    \Hom_{\S_n}\left(\sgn, R_n^{(k,\ell)} \otimes \bigwedge\nolimits^d \C^{n-1}\right) \\
    &\cong 
    \Hom_{\S_n}\left(\bigwedge\nolimits^{n-d-1}\C^{n-1}, R_n^{(k,\ell)}\right). 
\end{aligned}
\end{equation}
Consequently, proving that $f\cdot p_{\r;\s;1}=0$ for all $f \in \left(R_n^{(k,\ell)} \otimes \bigwedge\nolimits^{d-1} \C^{n-1}\right)^{\sgn}$ is equivalent to showing that $\Psi_{f\cdot p_{\r;\s;1}}=0$ in $\Hom_{\S_n}\left(\bigwedge\nolimits^{n-d-1}\C^{n-1}, R_n^{(k,\ell)}\right)$.

To prove $\Psi_{f\cdot p_{\r;\s;1}}=0$, in Theorem~\ref{thm:spanning} we construct a spanning set 
for $\Hom_{\S_n}\left(\bigwedge\nolimits^{n-d}\C^{n}, R_n^{(k,\ell)}\right)$ of antisymmetrization operators 
$\psi_{\m}$ indexed by monomials $\m$. 
In particular, these operators satisfy the special property,
that when multiplied by certain single index $i$ monomials and then summed over all $1\leq i \leq n$, 
the sum vanishes, see Lemma~\ref{lem:Zero-sum}.
Since the operators $\Psi_{f\cdot p_{\r;\s;1}}$ can be expressed as 
a linear combination of products of the spanning set $\psi_{\m}$ with certain single index monomials, 
then $\Psi_{f\cdot p_{\r;\s;1}}=0$ by the above property, so $\pi^{\sgn}$ is injective. 
Since $\pi$ is always a surjective map, the restriction to the sign-isotypic component is then an isomorphism.
Our main result, proven in Theorem~\ref{thm:iso-sign}, is the following. 

\begin{introthm}\label{thm:INTRO iso-sign}
    For any $n\geq 1$ and $k, \ell \geq 0$, the multiplication map $\pi: R_n^{(k,\ell)} \otimes R_n^{(0,1)} \to R_n^{(k,\ell+1)}$ is an isomorphism on sign-isotypic components. 
\end{introthm}

As an immediate consequence, from equation~\eqref{eq:INTRO-hook-hom-iso}, 
we obtain that the sign character of $R_n^{(k,\ell+1)}$ equals 
the sum of the hook characters of $R_n^{(k,\ell)}$, see Corollary~\ref{cor:hook-to-sign-sf}.

\begin{introcor}\label{cor:INTRO hook-to-sign-sf}
Let $n\geq 1$, let $k,\ell \geq 0$, and let $a$ record the degree in the alphabet $\bm{\theta}^{(\ell+1)}$. 
    Then
    \begin{align*}
        \left\langle \Frob(R_n^{(k,\ell+1)}; \mathbf{q};\mathbf{u}, a), s_{(1^n)} \right\rangle
   &= \sum_{d=0}^{n-1} a^d \left\langle \Frob(R_n^{(k,\ell)}; \mathbf{q};\mathbf{u}), s_{(d+1,1^{n-d-1})}\right\rangle.
    \end{align*}
\end{introcor}

The results and constructions above extend mutatis mutandis to 
the bosonic-fermionic coinvariant rings for the hyperoctahedral group $\mathfrak{B}_n$, 
with only minor modifications, see Section~\ref{sec:type-B}. 
In particular, Theorem \ref{thm:iso-det} and Corollary~\ref{cor:hook-to-sign-sf-B} provide the appropriate type B analogues to Theorem~\ref{thm:INTRO iso-sign} and Corollary~\ref{cor:INTRO hook-to-sign-sf}. 

\subsection{The socle of $R_n^{(k,\ell)}$} One of Haiman's seminal results on diagonal coinvariants is the 
\emph{operator theorem}, conjectured initially in \cite{Haiman1994} and proven geometrically in \cite{Haiman2002}. 
Stated for the diagonal harmonics $H_n^{(2,0)}$, this theorem asserts that 
the action of certain polarization operators on the classical Vandermonde determinant, 
along with certain partial differential operators $\partial_{x_i}$, generate $H_n^{(2,0)}$. 
A consequence of this theorem is that $H_n^{(2,0)} = \C[\partial_{x_1},\ldots,\partial_{x_n}]\left(H_n^{(2,0)}\right)^{\sgn}$, 
which when dualized
states that the socle of $R_n^{(2,0)}$ coincides with its sign-isotypic component $\left(R_n^{(2,0)}\right)^{\sgn}$.

As a particularly nice consequence of Lemma~\ref{lem:Zero-sum}, in Theorem~\ref{thm:socle} we obtain a generalization of one containment of Haiman's result to all bosonic-fermionic coinvariant rings.

\begin{introthm}\label{thm:INTRO socle} For all $n\geq 1$ and $k,\ell \geq 0$ we have that
    \[\left(R_n^{(k,\ell)}\right)^{\sgn} \subseteq \soc\left(R_n^{(k,\ell)}\right).\]
\end{introthm}

\subsection{Beyond $(0,1)$ and other isotypic components} In Sections \ref{sec:miniswaddle} and \ref{sec:mezzoswaddle} we prove that Theorem \ref{thm:INTRO iso-sign} fails to hold for other isotypic components or when tensoring with arbitrary coinvariant rings $R_n^{(a,b)}$. 

In particular, in Theorem~\ref{thm:neg-result} we show that for $n\geq 3$ and nontrivial values  
of $(k,\ell)$ and $(a,b)$ the surjection 
$\pi: R_{n}^{(k,\ell)} \otimes R_{n}^{(a,b)} \twoheadrightarrow R_{n}^{(k+a,\ell+b)}$ 
fails to be an isomorphism unless one of them is $(0,1)$. 

\begin{introthm}\label{thm:INTRO neg-result}
Let $n \geq 3$.
The surjection
\begin{equation*}
        \pi: R_{n}^{(k,\ell)} \otimes R_{n}^{(a,b)} \twoheadrightarrow R_{n}^{(k+a,\ell+b)}
    \end{equation*}
     is an isomorphism on sign-isotypic components if and only if 
     at least one of $(k,\ell)$ or $(a,b)$ is $(0,0)$ or $(0,1)$. 
\end{introthm}

Since our coinvariant rings admit natural inclusions $R_n^{(a,b)} \subseteq R_n^{(a',b')}$ whenever 
$a\leq a'$ and $b\leq b'$, Theorem \ref{thm:INTRO neg-result} reduces to showing 
$\pi: R_{n}^{(k,\ell)} \otimes R_{n}^{(a,b)} \twoheadrightarrow R_{n}^{(k+a,\ell+b)}$ 
is not an isomorphism whenever $(k,\ell), (a,b)\in \{(1,0), (0,2)\}$, 
which is the content of Propositions \ref{prop:(1,0)x(1,0)}, \ref{prop:(1,0)x(0,2)}, and \ref{prop:(0,2)x(0,2)}. 
When $n \leq 2$ the map degenerates to isomorphism for any value of 
$(k,\ell)$ and $(a,b)$, see Proposition \ref{prop:small-n-mini}. 

To conclude, we also study the restriction of $\pi: R_{n}^{(k,\ell)} \otimes R_{n}^{(0,1)} \twoheadrightarrow R_{n}^{(k,\ell+1)}$ to other isotypic components. 
The kernel of $\pi$ always contains a copy of the maximal ideal of $R_n^{(k,\ell)}$ by Lemma \ref{lem:K_{n-1}}, so for $k \geq 1$ the restriction of $\pi$ is never injective outside of the sign-isotypic component, see Theorem \ref{thm:mezzo-swaddle}.

\begin{introthm}\label{thm:INTRO mezzo-swaddle}
    Fix $n \geq 1$, $k \geq 1$, and $\ell \geq 0$, and let $\nu \vdash n$.
    Then the map
    \[
        \pi^\nu: \left(R_n^{(k,\ell)} \otimes R_n^{(0,1)}\right)^\nu \twoheadrightarrow \left(R_n^{(k,\ell+1)}\right)^\nu
    \]
    is injective if and only if $\nu = (1^n)$.
\end{introthm}

\subsection{Combinatorial consequences}
Using Corollary~\ref{cor:INTRO hook-to-sign-sf} 
we can determine the sign multiplicity of $R_n^{(k,\ell+1)}$ 
from the hook multiplicities of $R_n^{(k,\ell)}$, and similarly in type $B$.
The bigraded Frobenius series of $R_{\Z_r \wr \S_n}^{(1,1)}$ were computed in \cite{BCJLR}, 
which we use in Corollaries \ref{cor:12-sign} and \ref{cor:12-sign-B} 
to record the graded Hilbert series of $\left(R_n^{(1,2)}\right)^{\sgn}$ and $\left(R_{\mathfrak{B}_n}^{(1,2)}\right)^{\det}$. 

\begin{introcor}\label{cor:INTRO 12-sign}
For $n \geq 1$,
    \begin{align*}
        \left\langle \Frob (R_n^{(1,2)}; q; u, v), s_{(1^n)} \right\rangle &= \sum_{0 \leq i+j \leq n-1} u^i v^j q^{\binom{n-i-j}{2}} \qbinom{n-1}{j}_q\qbinom{n-1-j}{i}_q, 
    \\
        \left\langle \Frob (R_{\mathfrak{B}_n}^{(1,2)}; q; u, v), s_{(\varnothing,(1^n))} \right\rangle &= \sum_{d=0}^n v^d \qbinom{n}{d}_{q^2} \prod_{j=1}^{n-d} (u+q^{2j-1}).
    \end{align*}
Consequently, the multiplicity of the sign character of $R_n^{(1,2)}$ is $3^{n-1}$ and the multiplicity of the determinant character of $R_{\mathfrak{B}_n}^{(1,2)}$ is $3^{n}$.
\end{introcor}

Corollary~\ref{cor:INTRO hook-to-sign-sf} also recovers the known sign characters of $R_n^{(1,1)}$, $R_n^{(0,2)}$, $R_n^{(2,1)}$ and $R_n^{(0,3)}$ and the determinant characters of $R_{\mathfrak{B}_n}^{(1,1)}$, $R_{\mathfrak{B}_n}^{(0,2)}$ and $R_{\mathfrak{B}_n}^{(0,3)}$, see Section~\ref{sec:consequences}.

Another combinatorial consequence of our work is a proof of 
F.~Bergeron's \emph{hook component conjecture} \cite{BergeronOPAC}*{Conj. 2}. 
By the diagonal supersymmetry theorem of the second 
author {\cite{Lentfer-Supersymmetry}*{Cor. 1.3}}, 
there exist universal series coefficients $c_{\lambda,\mu} \in \Z_{\geq0}$ 
independent of $(k,\ell)$, indexed by partitions $\lambda$ with 
at most $n$ parts and $\mu \vdash n$, such that 
\[
    \Frob(R_n^{(k,\ell)}; \mathbf{q}; \mathbf{u}) = \sum_{\lambda,\mu} c_{\lambda,\mu} s_\lambda(\mathbf{q}/\mathbf{u}) s_\mu
\]
where $s_\lambda(\mathbf{q}/\mathbf{u})$ is a super Schur function evaluated in the grading variables (see Theorem~\ref{thm:DSS-S_n}).
Bergeron conjectured that for all $0\leq d\leq n-1$ and all partitions $\rho$,
\[ c_{\rho,(d+1,1^{n-d-1})} = \sum_{\lambda} c_{\lambda,(1^n)}, \]
where the sum is over all partitions $\lambda\supseteq\rho$ 
such that $\lambda/\rho$ is a vertical strip of size $d$.
We prove this conjecture in Corollary \ref{cor:hook-components} 
and its type B analog in Corollary \ref{cor:hook-components-B}. 

\subsection{Outline of the paper} 
In Section~\ref{sec:background} we review the necessary background 
on $\S_n$-modules, bosonic-fermionic coinvariant rings, 
and define several maps and constructions used throughout. 
Section~\ref{sec:main} contains the statement and proof of 
our flagship result Theorem~\ref{thm:INTRO iso-sign}; subsequently we prove 
Theorem~\ref{thm:INTRO socle}. In Sections~\ref{sec:miniswaddle} 
and~\ref{sec:mezzoswaddle}, we prove Theorems~\ref{thm:INTRO neg-result} 
and~\ref{thm:INTRO mezzo-swaddle}, respectively. We then extend 
Theorem~\ref{thm:INTRO iso-sign} to type B in Section~\ref{sec:type-B}, 
and finally discuss the character formulas of Corollary~\ref{cor:INTRO 12-sign} 
and the proof of Bergeron's hook component conjecture in Section~\ref{sec:consequences}.

%%%%%%%%%%%%%%%%%%%%%%%%%%%%%%%%%%%%%%%%%%%%%%%%%%%%%%%%%%%%%%%%%%%%%%
\section*{Acknowledgements}
%%%%%%%%%%%%%%%%%%%%%%%%%%%%%%%%%%%%%%%%%%%%%%%%%%%%%%%%%%%%%%%%%%%%%%

We thank Eugene Gorsky, Brendon Rhoades, and Josh Swanson for helpful conversations. 
N.G. was supported by the Natural Sciences and 
Engineering Research Council of Canada, RGPIN-2026-07206. 
P.S.O. was partially supported by JSPS KAKENHI grant JP25K17242. 

%%%%%%%%%%%%%%%%%%%%%%%%%%%%%%%%%%%%%%%%%%%%%%%%%%%%%%%%%%%%%%%%%%%%%%
\section{Preliminaries on bosonic-fermionic coinvariant rings}\label{sec:background}
%%%%%%%%%%%%%%%%%%%%%%%%%%%%%%%%%%%%%%%%%%%%%%%%%%%%%%%%%%%%%%%%%%%%%%

Let $n \in \mathbb{N}_{>0}$ be a fixed natural number. We denote by 
$\bm{z} \coloneqq \{ z_1, \ldots, z_n \}$ a family of $n$ variables, which 
will be called an \newword{alphabet}. 
Let $k, \ell \geq 0$. We denote by $\bm{x}^{(1)}, \ldots, \bm{x}^{(k)}$ 
a collection of $k$ commutative, also called bosonic, alphabets. 
We denote by $\bm{\theta}^{(1)}, \ldots, \bm{\theta}^{(\ell)}$ 
a collection of $\ell$ anticommutative, also called fermionic, alphabets. 
Bosonic and fermionic alphabets commute with each other, so that for any 
$\bullet$ and $\star$ we have the relations 
\[
    \left[x^{(\bullet)}_i,x^{(\star)}_j\right] 
    = 0 
    = \left[x^{(\bullet)}_i,\theta^{(\star)}_j\right], 
    \quad 
    \theta^{(\bullet)}_i\theta^{(\star)}_j + \theta^{(\star)}_j\theta^{(\bullet)}_i =0, 
    \quad 
    \text{and} 
    \quad 
    \left(\theta^{(\bullet)}_i\right)^2 =0. 
\]
Thus the vector space 
\begin{align*}
    \C\left[\bm{x}^{(1)}, \ldots, \bm{x}^{(k)}, 
    \bm{\theta}^{(1)}, \ldots, \bm{\theta}^{(\ell)}\right] 
    &= \C\left[\bm{x}^{(1)}, \ldots, \bm{x}^{(k)}\right] \otimes 
    \bigwedge\left[\bm{\theta}^{(1)}, \ldots, \bm{\theta}^{(\ell)}\right]\\  
    &\cong \Sym\left(\left(\C^n\right)^{\oplus k}\right) \otimes 
    \bigwedge \left(\left(\C^n\right)^{\oplus \ell}\right) 
\end{align*}
is an $\mathfrak{S}_n$-module via the \newword{diagonal action} given by 
permuting letters within each alphabet simultaneously. 
Namely, given $\sigma \in \mathfrak{S}_n$, then for all $1 \leq i \leq n$ and 
all $\bullet$ we have 
\[
    \sigma(x^{(\bullet)}_i) = x^{(\bullet)}_{\sigma(i)} 
    \quad \text{and} \quad 
    \sigma(\theta^{(\bullet)}_i) = \theta^{(\bullet)}_{\sigma(i)}. 
\]

The \newword{$(k,\ell)$-bosonic-fermionic coinvariant ring} is the quotient 
\begin{equation}
    \label{eq:coinv} 
    R_n^{(k,\ell)} \coloneqq 
    \frac{\C\left[\bm{x}^{(1)}, \ldots, \bm{x}^{(k)}, \bm{\theta}^{(1)}, \ldots, 
    \bm{\theta}^{(\ell)}\right]}{ \left\langle 
    \C\left[\bm{x}^{(1)}, \ldots, \bm{x}^{(k)}, 
    \bm{\theta}^{(1)}, \ldots, \bm{\theta}^{(\ell)}\right]^{\mathfrak{S}_n}_+ \right\rangle },
\end{equation}
whose defining ideal is generated by the polynomials having no constant 
term and invariant under the diagonal action of $\mathfrak{S}_n$, namely 
\[
 I_n^{(k,\ell)} \coloneqq 
\left\langle \left[\bm{x}^{(1)}, \ldots, \bm{x}^{(k)}, \bm{\theta}^{(1)}, \ldots, 
\bm{\theta}^{(\ell)}\right]^{\mathfrak{S}_n}_+
\right\rangle.
\]
We abbreviate $\C_n^{(k,\ell)} \coloneqq 
\C\left[\bm{x}^{(1)},\ldots,\bm{x}^{(k)},\bm{\theta}^{(1)},\ldots,\bm{\theta}^{(\ell)}\right]$, 
so that $R_n^{(k,\ell)} = \C_n^{(k,\ell)} / I_n^{(k,\ell)}$. 

It is sometimes convenient to consider the spaces above restricted to certain alphabets. 
Let $[n] \coloneqq \{1,\ldots, n\}$ and consider any nonempty subset $\alpha \subseteq [n]$. 
We denote by
$\bm{z}_\alpha\coloneqq\{z_i\mid i \in \alpha\}$ the subset of alphabet $\bm{z}$ restricted to $\alpha$, and set
\[
    \C_\alpha^{(k,\ell)} \coloneqq 
    \C\left[\bm{x}_\alpha^{(1)},\ldots,\bm{x}_\alpha^{(k)}, 
    \bm{\theta}_\alpha^{(1)},\ldots,\bm{\theta}_\alpha^{(\ell)}\right]
    \qquad \text{and}\qquad
    I_\alpha^{(k,\ell)} \coloneqq 
    \left\langle\left(\C_\alpha^{(k,\ell)}\right)^{\S_\alpha}_+\right\rangle.
\]

Then, for any nonempty subset $\alpha \subseteq [n]$, and any sequences $\r = (r_1,\ldots,r_k) \in (\Z_{\geq 0})^k$ and $\s= (s_1,\ldots,s_\ell) \in \{0,1\}^\ell$, we define the 
\newword{(super) polarized power-sum} symmetric function 
$p_{\r;\s}(\bm{x}^{(\bullet)}_\alpha;\bm\theta^{(\star)}_\alpha)$ as the sum 
\[
    p_{\r;\s}\left(\bm{x}^{(\bullet)}_\alpha;\bm\theta^{(\star)}_\alpha\right) 
    \coloneqq 
    \sum_{i \in \alpha} \prod_{j=1}^k\left(x_i^{(j)}\right)^{r_j} 
    \prod_{j=1}^\ell\left(\theta_i^{(j)}\right)^{s_j} 
    = \sum_{i \in \alpha}  m_{\r;\s}(i), 
\]
where
\[
    m_{\r;\s}(i) = 
    \prod_{j=1}^k\left(x_i^{(j)}\right)^{r_j} 
    \prod_{j=1}^\ell\left(\theta_i^{(j)}\right)^{s_j} 
    = 
    \left(x_i^{(1)}\right)^{r_1} \cdots \left(x_i^{(k)}\right)^{r_k} 
    \left(\theta_i^{(1)}\right)^{s_1} \cdots \left(\theta_i^{(\ell)}\right)^{s_\ell}. 
\]
We will often abbreviate and write 
$p_{\r;\s}(\alpha)\coloneqq  
p_{\r;\s}\left(\bm{x}^{(\bullet)}_\alpha;\bm\theta^{(\star)}_\alpha\right)$ 
and, when $\alpha = [n]$, omit the label entirely 
by writing $p_{\r;\s}\coloneqq p_{\r;\s}([n])$.

For any sequence $\s=(s_1,\ldots,s_m)$ let $|\s|\coloneqq \sum_{i=1}^m s_i$. 
We recall Weyl's polarization theorem.
\begin{proposition}[\cite{Weyl}; this formulation \cite{JiangLentfer}*{Prop. 2.1}]
\label{prop:ideal-generators}
    The (nonunital) algebra $(\C^{(k,\ell)}_\alpha)_+^{\S_\alpha}$, and thus the ideal
    $I_\alpha^{(k,\ell)}$,
    is generated by the polarized power sums 
    \begin{equation}\label{eq:gen-set-A-21}
        \left\{ p_{\r;\s}(\alpha) \mid
        \r \in (\Z_{\geq 0})^k,\ \s \in \{0,1\}^\ell, 1 \leq |\r|+|\s| \leq |\alpha| \right\}.
\end{equation}
    In particular, the set~\eqref{eq:gen-set-A-21} at $\alpha = [n]$ 
    generates the defining ideal $I_{n}^{(k,\ell)}$. 
\end{proposition}

The algebra $\C_n^{(k,\ell)}$ has a natural multigrading by alphabets, by setting 
\begin{align*}
    \deg \left(
    \prod_{j = 1}^{k} \prod_{i = 1}^{n} \left(x_i^{(j)}\right)^{r_{ij}} 
    \prod_{j = 1}^{\ell} \prod_{i = 1}^{n} \left(\theta_i^{(j)}\right)^{s_{ij}} 
    \right) = \left( 
    \sum_{i = 1}^{n} r_{i1}, \ldots, \sum_{i = 1}^{n} r_{ik}, 
    \sum_{i = 1}^{n} s_{i1}, \ldots, \sum_{i = 1}^{n} s_{i\ell}
    \right)&\\ 
    \in \left(\mathbb{Z}_{\geq0}\right)^{k} \times \{0,\ldots,n\}^{\ell},& 
\end{align*}
since a monomial has degree at most $n$ in each fermionic alphabet $\bm{\theta}$.
Then the algebras $R_n^{(k,\ell)}$ inherit this multigrading. 
In fact, $R_n^{(k,\ell)}$ vanishes in degree $n$ in each fermionic alphabet, 
because $\theta_n^{(j)} \equiv -\sum_{i<n}\theta_i^{(j)} \pmod{I_n^{(k,\ell)}}$ implies 
$\theta_1^{(j)}\cdots\theta_n^{(j)} \equiv 0 \pmod{I_n^{(k,\ell)}}$. Hence, $R_n^{(k,\ell)}$ is a
$\left(\mathbb{Z}_{\geq 0}\right)^{k} \times \{0,\ldots,n-1\}^{\ell}$-graded 
$\mathfrak{S}_n$-module, and decomposes as
\[
    R_n^{(k,\ell)}
    = \bigoplus_{(\r, \s) 
    \in (\mathbb{Z}_{\geq 0})^{k} \times \{0,\ldots,n-1\}^{\ell}} 
    \left(R_n^{(k,\ell)}\right)_{(\r, \s)},
\]
where each of the multihomogeneous components $\left(R_n^{(k,\ell)}\right)_{(\r,\s)}$ is an $\mathfrak{S}_n$-module. 

It is not difficult to see that the degree components of $R_n^{(k,\ell)}$ are invariant under permutation.
Let $\sigma(\r)$ denote the sequence with $\sigma(\r)_{\sigma(j)} = r_j$, and similarly let $\tau(\s)$ denote the sequence with $\tau(\s)_{\tau(j)} = s_j$.

\begin{lemma}
    \label{lem:degree-permutation}
    Let $k, \ell \geq 0$, 
    let  $\r=(r_1,\ldots, r_k) \in \mathbb{Z}_{\geq 0}^k$, 
    let $\s=(s_1,\ldots,s_\ell) \in \{0,\ldots,n-1\}^{\ell}$, 
    and let $(\sigma,\tau) \in \S_k \times \S_\ell$. 
    Then there is an isomorphism on 
    the multigraded components that is equivariant with respect to the diagonal 
    $\S_n$-action, namely 
    \begin{equation}
        \label{eq:permute graded components}
        \left(R_n^{(k,\ell)}\right)_{(\r,\s)} \cong 
        \left(R_n^{(k,\ell)}\right)_{(\sigma(\r),\tau(\s))}.
    \end{equation}
\end{lemma}

\begin{proof}
    Consider the action of $\S_k \times \S_\ell$ on $\C_n^{(k,\ell)}$ 
    defined on each variable as 
    \[
         (\sigma, \tau): x_i^{(j)} \mapsto x_i^{(\sigma(j))} 
         \qquad \text{and} \qquad 
         (\sigma, \tau): \theta_i^{(j)} \mapsto \theta_i^{(\tau(j))}, 
    \]
    for all $(\sigma, \tau) \in \S_k \times \S_\ell$, 
    and extended multiplicatively. 
    This action commutes with the diagonal $\S_n$ action by construction, 
    and we now show that it descends to an action of 
    $\S_k \times \S_\ell$ on $R_n^{(k,\ell)}$. 
    The map is well defined on $R_n^{(k,\ell)}$ because $(\sigma, \tau) \in \S_k \times \S_\ell$ maps $\left(\C_n^{(k,\ell)}\right)^{\S_n}_+$ onto itself.
    If $g \in \left(\C_n^{(k,\ell)}\right)^{\S_n}_+$, then
    for all $w\in \S_n$ we have 
    \[
        w \circ (\sigma,\tau)(g) = 
        (\sigma,\tau) \circ w(g) = 
        (\sigma,\tau)  (g), 
    \]
    so $(\sigma,\tau)(g)$ is $\S_n$-invariant. Since $(\sigma,\tau)$ is degree-preserving, $(\sigma,\tau)(g)$ has no constant term.
    Applying the same argument to $(\sigma,\tau)^{-1}$ gives $(\sigma,\tau)\left(\C_n^{(k,\ell)}\right)^{\S_n}_+ = \left(\C_n^{(k,\ell)}\right)^{\S_n}_+$,
    so $(\sigma,\tau) I_n^{(k,\ell)} = I_n^{(k,\ell)}$. 
    Moreover, given $(\sigma,\tau) \in \S_k \times \S_\ell$ then 
    for all $f, g \in \C_n^{(k,\ell)}$ we have 
    $(\sigma,\tau)(f) (\sigma,\tau)(g) = (\sigma,\tau)(fg)$ so 
    $(\sigma,\tau)(\theta_i^{(j)}) (\sigma,\tau)(\theta_i^{(j)}) = 
    (\sigma,\tau)(\theta_i^{(j)}\theta_i^{(j)}) = 0$, so indeed 
    $(\sigma, \tau) \in \mathrm{Aut}_{\S_n}\left(R_n^{(k,\ell)}\right)$ 
    and the map gives an action on $R_n^{(k,\ell)}$. 

    Consequently, any $(\sigma,\tau)$ gives an isomorphism on $R_n^{(k,\ell)}$, 
    and it clearly permutes the graded components as claimed 
    in equation~\eqref{eq:permute graded components}. 
\end{proof}

\subsection{Frobenius characters and symmetric functions}

We recall some basic facts but refer the reader to standard sources such as \cite{Macdonald}*{Chap. 1}, \cite{StanleyEC2}*{Chap. 7}, and \cite{fulton-youngtableaux}*{Chap. 7} for a thorough treatment on the representations of $\S_n$ and the ring of symmetric functions. 

Recall that the irreducible representations of $\C[\S_n]$-mod are given by the \newword{Specht modules} $S^\lambda$, which are indexed by partitions $\lambda\vdash n$. These representations provide a complete orthogonal set of simple modules for this category. Given any $\S_n$-module $M$, we denote by $M^\lambda$ its $\lambda$-isotypic component and by $M_\lambda$ its corresponding multiplicity space, so that
\[
    M \cong \bigoplus_{\lambda \vdash n} M^\lambda \cong 
    \bigoplus_{\lambda \vdash n} S^\lambda \otimes M_\lambda 
    \qquad \text{and} \qquad
    M_\lambda \cong \Hom_{\S_n}(S^\lambda, M).
\]
In the case when $\lambda$ is a column, we denote the sign representation by $\sgn = S^{(1^n)}$.

Let $\mathbf{q} = (q_1,\ldots, q_k)$ and $\mathbf{u} = (u_1,\ldots, u_\ell)$ be families of indeterminates.
The \newword{ring of symmetric functions} $\Lambda_{\mathbf{q}, \mathbf{u}}$ in infinitely many variables with coefficients in $\C[\mathbf{q},\mathbf{u}]$ is graded by homogeneous degree $\Lambda_{\mathbf{q}, \mathbf{u}} = \bigoplus_{n \in \Z_{\geq0}} \Lambda_{\mathbf{q}, \mathbf{u}}^n$, where each homogeneous component $\Lambda_{\mathbf{q}, \mathbf{u}}^n$
has a $\Z[\mathbf{q},\mathbf{u}]$-basis given by the \newword{Schur functions} $s_\lambda$ as $\lambda$ varies over all partitions $\lambda \vdash n$. 

The category $\C[\S_n]$-gmod of $(k+\ell)$-graded $\S_n$-modules categorifies the ring $\Lambda_{\mathbf{q}, \mathbf{u}}^n$ via the \newword{Frobenius character map}.
This map assigns to each module $M \in \C[\S_n]$-gmod a symmetric function in infinitely many variables, such that to each Specht module $\Frob(S^\lambda) = s_\lambda$. In particular, the Hom pairing on $\C[\S_n]$-gmod descends to the 
\newword{Hall inner product} on $\Lambda_{\mathbf{q}, \mathbf{u}}^n$, where the orthogonality of the Specht modules corresponds to the orthonormality of the Schur functions.

\begin{definition}
    \label{defi:multigraded-hilbert-frobenius-series}
    The \newword{multigraded Frobenius series} of $R_n^{(k,\ell)}$ is 
    \[
        \Frob(R_n^{(k,\ell)}; \mathbf{q}; \mathbf{u}) \coloneqq 
        \sum_{(\r, \s) 
    \in (\mathbb{Z}_{\geq 0})^{k} \times \{0,\ldots,n-1\}^{\ell}}
        \Frob\left(R_n^{(k,\ell)}\right)_{(\r, \s)} 
        q_1^{r_1}\cdots q_k^{r_k} u_1^{s_1}\cdots u_\ell^{s_\ell}, 
    \]
    where $\mathbf{q} = (q_1, \ldots, q_k)$ and $\mathbf{u} = (u_1, \ldots, u_{\ell})$ are tuples of variables.
\end{definition}

\subsection{The exterior algebra \texorpdfstring{$R_n^{(0,1)}$}{Rn(0,1)}}\label{subsec:Rn(0,1)}

It will be convenient to realize $R_n^{(0,1)}$ as an exterior algebra with basis given as follows. 
Let $e_1,\ldots,e_{n}$ be the standard basis for $\C^{n}$, taken as an $\S_n$-module under the usual permutation action $\sigma e_i = e_{\sigma(i)}$, so that
\begin{equation} \label{eq:R(0,1)=hooks}
R_n^{(0,1)} \cong \bigwedge \C^n / \langle e_1+\cdots +e_n\rangle
\cong \bigwedge \C^{n-1} = \bigoplus_{d=0}^{n-1} \bigwedge\nolimits^d \C^{n-1} \cong \bigoplus_{d=0}^{n-1} S^{(n-d,1^d)}, 
\end{equation}
where $\C^{n-1} = \C^n /\langle e_1 + \cdots + e_n\rangle \cong S^{(n-1,1)}$. By continuing to write $e_1, \ldots, e_n$ for the images of the basis vectors in $\C^{n-1}$, with $e_n= -(e_1+ \cdots +e_{n-1})$, we have that $e_1,\ldots, e_{n-1}$ generate $R_n^{(0,1)}$. 

Since $R_n^{(0,1)}$ is singly graded with degree $d$ 
component given by $\left(R_n^{(0,1)}\right)_d \cong \bigwedge^d \C^{n-1} \cong S^{(n-d,1^d)}$, we define a basis for each graded piece as follows. 

For any subset $\tau = \{ \tau_1<\cdots<\tau_d\} \subset [n-1]$ 
with $|\tau|=d$, set 
$e_\tau \coloneqq e_{\tau_1} \wedge \cdots \wedge e_{\tau_d}$ and let $\overline{\tau}= [n-1]\setminus \tau$.  Then, the set $\{ e_{\tau} \;:\; |\tau|=d\}$ forms a basis for $\bigwedge^d\C^{n-1}$, and the set $\{ e_{\overline{\tau}} \;:\; |\tau|=d\}$ forms a basis for $\bigwedge^{n-1-d}\C^{n-1}$. 
In this presentation, $\S_n$ acts by $\sigma(e_\tau) = e_{\sigma(\tau)}\coloneqq e_{\sigma(\tau_1)} \wedge \cdots \wedge e_{\sigma(\tau_d)} $, where $e_n = -(e_1 + \cdots + e_{n-1})$ is substituted whenever $\sigma(\tau_a) = n$ for some $a \in [d]$. In particular, when $\sigma \in \S_\tau$ then $\sigma(e_\tau) = \sgn(\sigma) e_\tau$.

We will also consider the element $\omega \coloneqq e_1\wedge \cdots \wedge e_{n-1}$, which generates the sign representation 
$\C \omega \cong \bigwedge^{n-1}\C^{n-1} \cong  \sgn.$

Since for any partition $\lambda \vdash n$ with transpose $\lambda^t$, we have $\sgn \otimes S^\lambda \cong S^{\lambda^t}$, and any $\S_n$-module $M$ over $\C$ is self dual, then we obtain the isomorphisms: 
\begin{equation}\label{eq:hom-isomorphism-hook-shapes}
\begin{aligned}
    \Hom_{\S_n}\left(\sgn, M \otimes S^{(n-d,1^d)}\right) 
    &\cong 
    \Hom_{\S_n}\left(\sgn \otimes  S^{(n-d,1^d)}, M\right) \\
    &\cong 
    \Hom_{\S_n}\left(S^{(d+1,1^{n-d-1})}, M\right).
\end{aligned}
\end{equation}
Hence, computing sign components of $M \otimes R_n^{(0,1)}$ reduces to computing hook components of $M$, 
\begin{equation}\label{eq:sign-hook}
\begin{aligned}
    \Hom_{\S_n}\left(\sgn, M \otimes R_n^{(0,1)}\right) 
    &\cong 
    \bigoplus_{d=0}^{n-1} \Hom_{\S_n}\left(S^{(d+1,1^{n-d-1})}, M\right) \\
    &\cong 
    \bigoplus_{d=0}^{n-1} \Hom_{\S_n}\left(\bigwedge\nolimits^{n-d-1}\C^{n-1}, M\right). 
\end{aligned}
\end{equation}
This identification, made precise in Proposition \ref{prop:Hom-iso}, is crucial for our main result in Section \ref{sec:main}.  

\subsection{Inclusions and projections of coinvariant rings} 
In this paper we will be considering various projections and inclusions of the aforementioned rings. We define them and state some of their properties below.

From the definitions we have the canonical quotient map of $\S_n$-modules, 
\begin{equation}\label{eq:quotient-map}
\Phi: \C_n^{(k,\ell)} \to R_n^{(k,\ell)}.
\end{equation}
sending each $f$ to its class in the quotient.

The inclusion of polynomials $\C_n^{(k,\ell)} \subset \C_n^{(k+a,\ell+b)}$, mapping each alphabet $\bm{z}^{(j)}$ to itself,
induces canonical multigraded $\S_n$-equivariant ring maps
\begin{align}
\iota &: R_n^{(k,\ell)} \hookrightarrow R_n^{(k+a,\ell+b)}, & f &\mapsto f, \label{eq:ring-inclusion}\\
\rho &: R_n^{(k+a,\ell+b)} \twoheadrightarrow R_n^{(k,\ell)}, & f &\mapsto f|_{\bm{x}^{(k+1)},\ldots, \bm{x}^{(k+a)}, \bm{\theta}^{(\ell+1)}, \ldots, \bm{\theta}^{(\ell+b)}=0}.
\end{align}
Since $\rho \circ \iota = \id$, we have that $\iota$ is an injection and $\rho$ is a surjection. 

The multiplication of polynomials 
\[
    \tilde{\pi} :
    \C_n^{(k,\ell)} \otimes \C_n^{(a,b)}\to \C_n^{(k+a,\ell+b)}
\]
is an isomorphism of multigraded (super)algebras which commutes with the diagonal 
action of $\S_n$ and preserves the gradings of $\C_n^{(k,\ell)}$ and 
$\C_n^{(a,b)}$ independently. 
We implicitly use the Koszul sign convention for superalgebras, namely for $f, f', g, g'$ 
homogeneous elements of (super)degree 
$\deg_{\Z_2}(f)$, $\deg_{\Z_2}(f')$, $\deg_{\Z_2}(g)$, $\deg_{\Z_2}(g')$ 
with respect to the (super)grading assigning $0$ to bosonic variables and 
$1$ to fermionic variables, their product is 
$(f \otimes g)(f' \otimes g') = (-1)^{\deg_{\Z_2}(g)\deg_{\Z_2}(f')}ff' \otimes gg'$. 

The following is standard.
\begin{lemma}[\cite{BourbakiAlgebraI}*{Cor. 1 p.~253}]\label{lem:tensor-quotients}
    Let $A$ and $A'$ be vector spaces with subspaces $I \subseteq A$ and $I' \subseteq A'$. Then the kernel of the surjection $A \otimes A' \twoheadrightarrow A/I \otimes A'/I'$ is $I \otimes A' + A \otimes I'$.
\end{lemma}

Under the multiplication map $\tilde{\pi}$, any $\S_n$-invariant polynomial of $\C_n^{(k,\ell)}$ or of $\C_n^{(a,b)}$ without constant term is also an  $\S_n$-invariant polynomial of $\C_n^{(k+a,\ell+b)}$ without constant term, so $I_n^{(k,\ell)} \otimes \C_n^{(a,b)} + \C_n^{(k,\ell)} \otimes I_n^{(a,b)} \subseteq I_n^{(k+a,\ell+b)}$.
Lemma~\ref{lem:tensor-quotients} then asserts that $\tilde{\pi}$ descends to a surjective $\S_n$-equivariant map, that preserves multidegrees, and whose kernel is the image of $I_n^{(k+a,\ell+b)}$ in $R_n^{(k,\ell)} \otimes R_n^{(a,b)}$:
\begin{align}\label{eq:pi-map}
    \pi : 
    R_n^{(k,\ell)} 
    \otimes 
    R_n^{(a,b)} 
    \twoheadrightarrow 
    R_n^{(k+a,\ell+b)}
    \qquad \text{sending} \qquad
    f\otimes g \mapsto fg.
\end{align}

We will often consider restrictions of these maps to isotypic components. For any $\S_n$-equivariant map $\phi: M \to N$ between any two $\S_n$-modules $M,N$ we will denote by 
\[
\phi^{\nu}: M^{\nu} \to N^{\nu}
\]
the restriction of $\phi$ to the corresponding $\nu$-isotypic components.

\begin{lemma}
    \label{lem:tensor-commutative-diagram}
    For any $k'\leq k, \ell'\leq \ell, a'\leq a, b'\leq b$, 
    there is a commutative, grading preserving diagram of $\S_n$-modules 
    \[
    \begin{tikzcd}
        R_n^{(k,\ell)} \otimes R_n^{(a,b)} \arrow[r, "\pi", two heads] & 
        R_n^{(k+a, \ell+b)} \\ 
        R_n^{(k',\ell')} \otimes R_n^{(a',b')} \arrow[r, "\pi'", two heads] \arrow[u, "\iota", hook] & 
        R_n^{(k'+a', \ell'+b')}\arrow[u, "\iota'", hook] 
    \end{tikzcd}.
    \]
    In particular, for any isotypic component $\nu$, if the restriction 
$\pi^{\nu}$ is injective then so is $(\pi')^{\nu}$. 
\end{lemma}

\begin{proof}
    Since for any $f \otimes g \in  R_n^{(k',\ell')} \otimes R_n^{(a',b')}$ 
    we have 
    \[
        \pi \circ \iota(f \otimes g) = 
        \pi(f \otimes g) = 
        fg = 
        {\iota'}(fg) = 
        {\iota'}\circ {\pi'}(f \otimes g), 
    \]
    the diagram commutes. 
    Moreover, since all four maps are $\S_n$-equivariant and grading preserving, 
    then so are their compositions. 

    Lastly, since $\iota'$ is an injection,  $\ker(\pi') = \ker(\iota' \circ \pi') = \ker(\pi \circ \iota) = \iota^{-1}(\ker \pi)$.
    Since $\iota$ is an injection and is $\S_n$-equivariant, it injects $(\ker \pi')^\nu$ into $(\ker \pi)^\nu$ for every $\nu$. Hence $(\pi')^\nu$ is injective whenever $\pi^\nu$ is.
\end{proof}

%%%%%%%%%%%%%%%%%%%%%%%%%%%%%%%%%%%%%%%%%%%%%%%%%%%%%%%%%%%%%%%%%%%%%%
%%%%%%%%%% MAIN %%%%%%%%%%%%%%%%%%%%%%%%%%%%%%%%%%%%%%%%%%%%%%%%%%%%%%
%%%%%%%%%%%%%%%%%%%%%%%%%%%%%%%%%%%%%%%%%%%%%%%%%%%%%%%%%%%%%%%%%%%%%%
\section{Sign components of \texorpdfstring{$R_n^{(k,\ell)} \otimes R_n^{(0,1)}$}{Rn(k,l)xRn(0,1)}}\label{sec:main}
%%%%%%%%%%%%%%%%%%%%%%%%%%%%%%%%%%%%%%%%%%%%%%%%%%%%%%%%%%%%%%%%%%%%%%

We now turn our attention to understanding the sign isotypic component of the ring $R_n^{(k,\ell+1)}$ for $k,\ell \geq 0$. To do so we specialize to $(a,b)=(0,1)$ in \eqref{eq:pi-map} and restrict our attention to the multiplication map 
\[
    \pi: R_n^{(k,\ell)} \otimes R_n^{(0,1)} \to R_n^{(k,\ell+1)}. 
\]
We will show that $\ker(\pi^{\sgn}) =0$, hence proving the sign-isotypic components
$\left(R_n^{(k,\ell)} \otimes R_n^{(0,1)}\right)^{\sgn}$ and $\left(R_n^{(k,\ell+1)}\right)^{\sgn}$ are isomorphic.

\subsection{Characterizing the kernel of \texorpdfstring{$\pi$}{pi}}\label{subsec:kernel-pi}
We begin by giving a precise characterization of $\ker(\pi)$.

Recall by Proposition~\ref{prop:ideal-generators} that the ideal $I_n^{(k,\ell+1)}$ is generated by polarized power sums $p_{\r;\s'}$ with $\r \in (\Z_{\geq 0})^k$ and $\s' \in \{0,1\}^{\ell+1}$. 
It will be convenient to explicitly factor out the $\bm{\theta}^{(\ell+1)}$ alphabet. 
Thus, letting $\varepsilon \in \{0,1\}$ denote the final entry in $\s'$ so that $\s' = (s_1,\ldots, s_\ell, \varepsilon)$ and $\s = (s_1,\ldots, s_\ell)$, we set $p_{\r;\s;\varepsilon}\coloneqq p_{\r;\s'}$. Hence, for $\bullet \in [k]$ and $\star \in [\ell]$ we have,
\begin{equation}
    p_{\r;\s;\varepsilon}\left(\bm{x}^{(\bullet)};\bm{\theta}^{(\star)};\bm{\theta}^{(\ell+1)}\right)  = \sum_{i=1}^n m_{\r;\s}(i) \left(\theta_i^{(\ell+1)}\right)^{\varepsilon},
\end{equation}
with the restrictions $p_{\r;\s;\varepsilon}(\alpha)$ and $p_{\r;\s;\varepsilon} = p_{\r;\s;\varepsilon}([n])$ interpreted as before.
In particular, we can restate Proposition \ref{prop:ideal-generators} as follows
\[
I_n^{(k,\ell+1)} = \left\langle p_{\r;\s;\varepsilon} \mid 
\r \in (\Z_{\geq 0})^k, \s \in \{0,1\}^\ell, \varepsilon \in \{ 0,1\}, \text{ with }
1 \leq |\r|+|\s| + \varepsilon \leq n \right\rangle.
\]

\begin{lemma}
    \label{lem:Ideal-decomposition}
    Let $K_n^{(k,\ell+1)} \coloneqq 
    \left\langle p_{\r;\s;1}
    \mid  \r \in (\mathbb{Z}_{\geq 0})^{k},\, 
    \s \in \{0,1\}^{\ell} \text{ with  } (\r,\s) \neq (\underline{0}, \underline{0}) \right \rangle \subset \C_n^{(k,\ell+1)}$. 
    Then
    \[
        I_n^{(k,\ell+1)} = 
        I_n^{(k,\ell)} \otimes \C_n^{(0,1)} + 
        \C_n^{(k,\ell)} \otimes I_n^{(0,1)} + 
        K_n^{(k,\ell+1)}. 
    \]
\end{lemma}

\begin{proof}
By Proposition~\ref{prop:ideal-generators}, $I_n^{(k,\ell)}\otimes \C_n^{(0,1)}$ is generated by the power sums $p_{\r;\s;0}$ and $\C_n^{(k,\ell)} \otimes I_n^{(0,1)}$ is generated by the power sums $p_{\underline{0};\underline{0};1}$. 
Thus, the only remaining generators are precisely the elements $p_{\r;\s;1}$ for which $(\r,\s ) \neq (\underline{0},\underline{0})$. 
\end{proof}

Intuitively, Lemma~\ref{lem:Ideal-decomposition} decomposes the defining ideal of 
$R_n^{(k,\ell+1)}$ into generators that only utilize the alphabets of $\C_n^{(k,\ell)}$ 
or only utilize the alphabets of $\C_n^{(0,1)}$, and precisely determines which generators 
utilize both alphabets. We can now explicitly compute the kernel of $\pi$ as follows. 

\begin{lemma}
    \label{lem:tensor-quotient-is-quotient-tensor}
    Let $\pi : R_n^{(k,\ell)} \otimes R_n^{(0,1)} \to R_n^{(k,\ell+1)}$ 
    be the surjective $\S_n$-equivariant map induced by the multiplication 
    of polynomials. Then 
    \begin{align*}
        \ker(\pi) &= 
        \left(R_n^{(k,\ell)} \otimes R_n^{(0,1)}\right) K_n^{(k,\ell+1)}\\
        &= 
       \left \langle f \cdot p_{\r;\s;1} 
        \mid   \r \in (\mathbb{Z}_{\geq 0})^{k},\, 
    \s \in \{0,1\}^{\ell} \text{ with  } (\r,\s) \neq (\underline{0}, \underline{0}),\, 
        f \in R_n^{(k,\ell)} \otimes R_n^{(0,1)} \right \rangle. 
    \end{align*}
\end{lemma}

\begin{proof}
Identifying the ideals $I_n^{(k,\ell+1)}$ and $K_n^{(k,\ell+1)}$ in $\C_n^{(k,\ell)} \otimes \C_n^{(0,1)} \cong \C_n^{(k,\ell+1)}$ with their images under $\Phi$ in $R_n^{(k,\ell)}\otimes R_n^{(0,1)}$, we have 
    \begin{align*}
        \ker(\pi) 
        & = (R_n^{(k,\ell)} \otimes R_n^{(0,1)}) I_n^{(k,\ell+1)} \\
        & = (R_n^{(k,\ell)} \otimes R_n^{(0,1)}) (I_n^{(k,\ell)} \otimes \C_n^{(0,1)}) 
         + (R_n^{(k,\ell)} \otimes R_n^{(0,1)}) (\C_n^{(k,\ell)} \otimes I_n^{(0,1)}) + 
        (R_n^{(k,\ell)} \otimes R_n^{(0,1)}) K_n^{(k,\ell+1)} \\ 
        & = (R_n^{(k,\ell)} \otimes \C_n^{(0,1)}) (I_n^{(k,\ell)} \otimes R_n^{(0,1)}) 
         + (\C_n^{(k,\ell)} \otimes R_n^{(0,1)}) (R_n^{(k,\ell)} \otimes I_n^{(0,1)}) + 
        (R_n^{(k,\ell)} \otimes R_n^{(0,1)}) K_n^{(k,\ell+1)} \\ 
        & = (R_n^{(k,\ell)} \otimes R_n^{(0,1)}) K_n^{(k,\ell+1)}. 
    \end{align*}
\end{proof}

As an immediate consequence, we obtain the following result. 

\begin{corollary}\label{cor:kernel-isotypic}
    Fix $n \geq 1$, and $k,\ell \geq 0$. For all $\nu \vdash n$ we have 
    \begin{align*}
        \ker(\pi)^{\nu} &= 
        \ker(\pi^{\nu}) = 
        \sum_{(\r,\s) \neq (\underline{0},\underline{0})} \left(R_n^{(k,\ell)} \otimes R_n^{(0,1)}\right)^{\nu} \cdot p_{\r;\s;1} \\
      &=  \left \langle f \cdot p_{\r;\s;1} \mid  \r \in (\mathbb{Z}_{\geq 0})^{k},\, \s \in \{0,1\}^{\ell} \text{ with  } (\r,\s) \neq (\underline{0}, \underline{0}),\, f \in \left(R_n^{(k,\ell)} \otimes R_n^{(0,1)}\right)^{\nu} \right \rangle. 
    \end{align*}
     In particular, $\pi^\nu$ is injective if and only if 
 $f \cdot p_{\r;\s;1} = 0$ for all $(\r,\s) \neq (\underline{0},\underline{0})$
 and all $f \in \left(R_n^{(k,\ell)} \otimes R_n^{(0,1)}\right)^{\nu}$. 
\end{corollary}

\subsection{Computing the alternating kernel}
With Corollary \ref{cor:kernel-isotypic} in hand, our attention now turns to understanding the sign isotypic component of the tensor product $R_n^{(k,\ell)}\otimes R_n^{(0,1)}$. Since $R_n^{(0,1)} \cong \bigwedge \C^{n-1}$, with generators $\{ e_\tau \mid \tau \subseteq [n-1]\}$, it is notationally convenient to write the map $\pi$ as follows:
\[
\pi: R_n^{(k,\ell)} \otimes \bigwedge \C^{n-1} \to R_n^{(k,\ell+1)}
\qquad \text{where} \qquad g \otimes e_i \mapsto g \theta_i^{(\ell+1)},
\]
with $e_n = -(e_1+\cdots+e_{n-1})$.

Now, recall from Section \ref{subsec:Rn(0,1)} that computing the sign component of $R_n^{(k,\ell)}\otimes R_n^{(0,1)}$ is equivalent to computing the hook isotypic components of $R_n^{(k,\ell)}$ (see equation~\eqref{eq:sign-hook}). 
In particular, since the degree $d$ component of $\bigwedge \C^{n-1}$ is given by $\bigwedge^d \C^{n-1}$ then it suffices to study the spaces $\Hom_{\S_n}\left(\bigwedge^{n-d-1}\C^{n-1}, M\right)$.

As before, for $0\leq d \leq n-1$ and $\tau = \{\tau_1<\cdots<\tau_d\} \subseteq [n-1]$ set $\overline{\tau}= [n-1]\setminus \tau$ and let $e_\tau = e_{\tau_1} \wedge \cdots \wedge e_{\tau_d}$ denote the basis elements of $\bigwedge^d \C^{n-1}$.  
Then, for each $0 \leq d \leq n-1$, define the bilinear pairing
 $\langle \cdot, \cdot \rangle: \bigwedge^d \C^{n-1} \times \bigwedge^{n-1-d} \C^{n-1} \to \C$ on the basis
\[
    (e_\tau, e_{\overline{\mu}}) \mapsto \langle e_\tau, e_{\overline{\mu}}\rangle 
\]
where $\langle e_\tau, e_{\overline{\mu}}\rangle =\pm \delta_{\tau,\mu}$ 
is the unique coefficient in the equation 
$e_\tau \wedge e_{\overline{\mu}} = \langle e_\tau, e_{\overline{\mu}}\rangle  \omega$ and $\omega = e_1 \wedge \cdots \wedge e_{n-1}$, so that more generally, $\eta \wedge \zeta = \langle \eta,\zeta \rangle \omega$. 

\begin{lemma}\label{lem:pairing-equivariant}
    Let $\C$ denote the $1$-dimensional sign representation.
    For any $0\leq d \leq n-1$,
    the map 
    $\langle \cdot, \cdot \rangle : \bigwedge^d \C^{n-1} \times \bigwedge^{n-1-d} \C^{n-1} \to\C $ 
    is a bilinear, $\S_n$-equivariant pairing, satisfying:
    \[
        \langle \sigma \eta, \sigma \zeta\rangle = \sgn(\sigma) \langle \eta,\zeta\rangle
    \]
    for any $\sigma \in \S_n$.
\end{lemma}

\begin{proof}
    The map is bilinear by construction. Moreover, since $\C \cong \C \omega = \sgn$ as an $\S_n$-module via the identification $1 \mapsto \omega $, then $\S_n$-equivariance follows since
    \[
    \langle \sigma \eta,\sigma \zeta \rangle \omega 
        = \sigma \eta \wedge\sigma \zeta
    = \sigma (\eta \wedge \zeta) 
        = \sgn(\sigma)( \eta\wedge \zeta)
        = \sgn(\sigma) \langle \eta, \zeta \rangle \omega
    \]
    hence, 
    $
\sigma(\eta,\zeta) 
        = (\sigma \eta,\sigma \zeta) 
        \mapsto \langle \sigma \eta,\sigma \zeta \rangle =
        \sgn(\sigma) \langle \eta, \zeta \rangle.
    $
\end{proof}

We now make the isomorphism in equation~\eqref{eq:sign-hook} relating the sign and hook isotypic components explicit. 

\begin{proposition} \label{prop:Hom-iso}Let $0 \leq d \leq n-1$.
    For any graded $\S_n$-module $M$,
    there is an isomorphism of graded vector spaces
    \[
        \Hom_{\S_n}\left( \sgn, M \otimes  \bigwedge\nolimits^d \C^{n-1}\right) 
        \to \Hom_{\S_n}\left( \bigwedge\nolimits^{n-d-1} \C^{n-1}, M\right)
    \]
    sending $F \mapsto \Psi_F$ where $F$ and $\Psi_F$ are the maps defined on the bases as follows:  
    \[F: 1 \mapsto \sum_{|\tau|=d} r_\tau \otimes e_\tau \qquad \text{and} \qquad
    \Psi_{F(1)}: e_{\overline{\mu}} \mapsto \sum_{ |\tau|=d}  \langle e_\tau , e_{\overline{\mu}}\rangle r_\tau = \langle e_\mu , e_{\overline{\mu}}\rangle r_\mu.
    \]
\end{proposition}

\begin{proof}
    We first prove that $\Hom_{\C}\left(\sgn, M \otimes \bigwedge\nolimits^d \C^{n-1}\right) \cong 
    \Hom_{\C}\left(\bigwedge\nolimits^{n-d-1} \C^{n-1}, M\right)$ via the map $F \mapsto \Psi_F$, 
    where we set $\Psi_F\coloneqq \Psi_{F(1)}$ for simplicity. 
    For injectivity, suppose $\Psi_F=0$. 
    Then $\Psi_F(e_{\overline{\mu}})=0$ for all $|\mu|=d$, so 
    \[
        0 = \Psi_F(e_{\overline{\mu}}) 
        = \sum_{|\tau| = d} \langle e_\tau, e_{\overline{\mu}}\rangle r_\tau = \pm r_\mu, 
    \]
    implying $r_\mu =0$ for all $|\mu|=d$, which is true if and only if 
    $F(1) = \sum_{|\tau| = d} r_\tau \otimes e_\tau =0$. 
    Thus $F=0$ and the map is injective. 
    For surjectivity, suppose $\Psi \in \Hom_{\C}\left(\bigwedge\nolimits^{n-d-1}\C^{n-1}, M\right)$. 
    Define $F_\Psi(1) = \sum_{|\tau|=d} \langle e_\tau, e_{\overline{\tau}} \rangle 
    \Psi(e_{\overline{\tau}}) \otimes e_\tau$, then for all $\beta$ with $|\beta|=d$ we have 
    \begin{equation}\label{eq:F-Psi-surjectivity}
        \Psi_{F_{\Psi}}(e_{\overline{\beta}}) = 
        \langle e_\beta, e_{\overline{\beta}} \rangle 
        \langle e_\beta, e_{\overline{\beta}} \rangle \Psi(e_{\overline{\beta}}) = 
        \Psi(e_{\overline{\beta}}). 
    \end{equation}
    Thus $\Psi_{F_\Psi} = \Psi$ and the assignment $F \mapsto \Psi_F$ is surjective. 

    Now we verify that the map $F \mapsto \Psi_F$ restricts to an isomorphism 
    $\Hom_{\S_n}\left(\sgn, M \otimes \bigwedge\nolimits^d \C^{n-1}\right) \cong 
    \Hom_{\S_n}\left(\bigwedge\nolimits^{n-d-1} \C^{n-1}, M\right)$. 
    Suppose $F \in \Hom_{\S_n}\left( \sgn, M \otimes \bigwedge\nolimits^d \C^{n-1}\right)$. 
    Then 
    \[
        F(1) 
        = \sgn(\sigma) F(\sigma(1)) 
        = \sgn(\sigma) \sigma F(1) 
        = \sgn(\sigma)\sum_{\tau} \sigma r_\tau \otimes \sigma e_\tau 
        = \sum_{\tau} (\sgn(\sigma) \sigma r_\tau) \otimes \sigma e_\tau, 
    \]
    so by Lemma \ref{lem:pairing-equivariant} we have 
    \[
        \Psi_F(\sigma e_{\overline{\mu}}) 
        = \sum_{|\tau| = d} \langle \sigma e_\tau, \sigma e_{\overline{\mu}} \rangle (\sgn(\sigma) \sigma r_\tau) 
        \overset{\ref{lem:pairing-equivariant}}{=} 
        \sum_{|\tau| = d} \langle e_\tau, e_{\overline{\mu}} \rangle \sigma r_\tau 
        = \sigma \left( \sum_{|\tau| = d} \langle e_\tau, e_{\overline{\mu}} \rangle r_\tau \right)
        = \sigma \Psi_F(e_{\overline{\mu}}) 
    \]
    hence $\Psi_F \in \Hom_{\S_n}\left( \bigwedge\nolimits^{n-d-1} \C^{n-1}, M\right)$. 
    Conversely, suppose $\Psi \in \Hom_{\S_n}\left(\bigwedge\nolimits^{n-d-1}\C^{n-1}, M\right)$. 
    Observe that for any 
    $G(1) = \sum_{|\tau|=d} r_\tau \otimes e_\tau \in M \otimes \bigwedge\nolimits^d\C^{n-1}$ we 
    have $\sigma G(1) = \sum_{|\tau|=d} \sigma r_\tau \otimes \sigma e_\tau$ for all $\sigma \in \S_n$. 
    Then for all $\beta$ with $|\beta|=d$ we have 
    \begin{align*}
        \Psi_{\sigma G(1)}(e_{\overline{\beta}}) 
        &= \sum_{|\tau| = d} \langle \sigma e_\tau, e_{\overline{\beta}} \rangle \sigma r_\tau 
        \overset{\ref{lem:pairing-equivariant}}{=} 
        \sum_{|\tau| = d} \sgn(\sigma) \langle e_\tau, \sigma^{-1} e_{\overline{\beta}} \rangle \sigma r_\tau \\
        &= \sgn(\sigma) \sigma \left(\sum_{|\tau| = d} 
        \langle e_\tau, \sigma^{-1} e_{\overline{\beta}} \rangle r_\tau \right) 
        = \sgn(\sigma) \sigma \Psi_{G(1)}( \sigma^{-1} e_{\overline{\beta}}), 
    \end{align*}
    so $\Psi_{\sigma G(1)} = \sgn(\sigma) \sigma \Psi_{G(1)} \sigma^{-1}$ 
    for all $G(1) \in M \otimes \bigwedge\nolimits^d \C^{n-1}$. 
    In particular, $G(1) = F_{\Psi}(1)$ gives 
    \[
        \Psi_{\sigma F_\Psi(1)} = 
        \sgn(\sigma) \sigma \Psi_{F_\Psi(1)}\sigma^{-1} = 
        \sgn(\sigma) \sigma \Psi \sigma^{-1} = 
        \sgn(\sigma) \Psi = 
        \sgn(\sigma) \Psi_{F_\Psi(1)} = 
        \Psi_{\sgn(\sigma) F_\Psi(1)}, 
    \]
    where we used the $\S_n$-equivariance of $\Psi$ in the third equality 
    and the linearity of the map $F \mapsto \Psi_F$ in the fifth equality. 
    Since the assignment $F \mapsto \Psi_F$ is injective, 
    we obtain $\sigma F_\Psi(1) = \sgn(\sigma) F_\Psi(1) = F_\Psi(\sigma 1)$, 
    therefore $F_{\Psi} \in \Hom_{\S_n}\left( \sgn, M \otimes \bigwedge\nolimits^d \C^{n-1}\right)$. 
\end{proof}

Recall that $\Hom_{\S_n}\left(\sgn, R_n^{(k,\ell)}\otimes \bigwedge^{d-1} \C^{n-1}\right) \cong \left(R_n^{(k,\ell)} \otimes \bigwedge^{d-1}\C^{n-1}\right)^{\sgn}$ via the map sending $F \mapsto F(1)$. Thus, since the sign component of $R_n^{(k,\ell)} \otimes R_n^{(0,1)}$ is graded by $R_n^{(0,1)}$, by equation~\eqref{eq:sign-hook} and Corollary \ref{cor:kernel-isotypic} it suffices to verify 
 that $F(1) \cdot p_{\r;\s;1} = 0$ for all $(\r,\s) \neq (\underline{0},\underline{0})$
 and all $F \in \Hom_{\S_n}\left(\sgn, R_n^{(k,\ell)}\otimes \bigwedge^{d-1} \C^{n-1}\right)$, for each $1\leq d \leq n$.
 Since by Proposition \ref{prop:Hom-iso}, $F(1) \cdot p_{\r;\s;1} = 0$ if and only if $\Psi_{F(1) \cdot p_{\r;\s;1}}$ is 0 in $\Hom_{\S_n}\left(\bigwedge^{n-d-1} \C^{n-1}, R_n^{(k,\ell)}\right)$, we proceed by proving the latter statement.

 To do so, we will build a spanning set for $\Hom_{\S_n}\left(\bigwedge^{n-d}\C^n, R_n^{(k,\ell)}\right)$ that vanishes when applied to any monomial (see Lemma~\ref{lem:Zero-sum}) and from it derive the identity above. 
 As in \cite{GLM}, it turns out to be simpler to construct a spanning set for the larger space of $(d,1^{n-d})$ and $(d+1,1^{n-d-1})$ isotypic components, and then restrict to the desired subpiece.
 Hence, with this in mind, we begin by identifying the isotypic components of $S^{(d,1^{n-d})} \oplus S^{(d+1,1^{n-d-1})}$ with the image of certain idempotent operators, whose construction we now recall.

For $\beta \subseteq [n]$, denote its Young subgroup in $\S_n$ by $\mathfrak{S}_\beta$. 
The \newword{antisymmetrizer} and 
\newword{symmetrizer} of $\mathfrak{S}_\beta$, respectively, are defined by 
\[
     \cY_\beta= \frac{1}{|\beta|!}\sum_{\sigma \in \mathfrak{S}_\beta} \sgn(\sigma) \sigma 
     \qquad \text{and} \qquad \cC_\beta= \frac{1}{|\beta|!}\sum_{\sigma \in \mathfrak{S}_\beta} \sigma. 
\]
Thus, an element  $f$ in an $\S_n$-module $M$ is \newword{$\S_\beta$-antisymmetric} or 
\newword{$\S_\beta$-symmetric}, respectively, whenever
$\cY_\beta(f) = f $ or $ \cC_\beta(f) = f.$
Then, for any $\S_n$-module $M$, we may consider the subset
\begin{align*}
    \cY_{\beta} \cC_{\overline{\beta}}(M)
    = 
    \left\{ f \in M : (\sigma,\tau)f = \sgn(\sigma) f 
    \text{ for all } 
    (\sigma,\tau) \in \S_{\beta} \times \S_{\overline{\beta}} \right\} 
    \subset M. 
\end{align*}

\begin{proposition}[\cite{GLM}*{Prop. 4.6}]
\label{prop:hook-hom-iso} Let $0\leq d \leq n-1$.
    For any $\mathfrak{S}_n$-module $M$ and $\beta \subseteq [n]$ with $|\beta|=n-d$, the assignment $\psi \mapsto \psi(e_\beta)$ for each $\psi \in \Hom_{\S_n}\left(\bigwedge\nolimits^{n-d}\C^n, M\right)$ induces an isomorphism
    \begin{align*}
\Hom_{\S_n}\left(\bigwedge\nolimits^{n-d}\C^n, M\right) \cong \cY_{\beta} \cC_{\overline{\beta}}(M).
    \end{align*}
\end{proposition}
Thus, we will construct a spanning set for $\Hom_{\S_n}\left(\bigwedge^{n-d}\C^n, R_n^{(k,\ell)}\right)$ by building one for $\cY_{\beta} \cC_{\overline{\beta}}\left(R_n^{(k,\ell)}\right)$ and pulling it back via the isomorphism. 
Our spanning set for $\cY_{\beta} \cC_{\overline{\beta}}\left(R_n^{(k,\ell)}\right)$ is  obtained by antisymmetrizing all possible monomials. 

\begin{definition}  Let $0\leq d \leq n-1$.
    Given monomials $m_i(j) = m_{\r^{(i)};\s^{(i)}}(j) \in R_n^{(k,\ell)}$ with $\r^{(i)},\s^{(i)}\geq \underline{0}$ for $1\leq i\leq n-d$, and any ordered $\tau \subseteq [n]$ with $\tau =(\tau_1,\dots,\tau_{n-d})$, let 
    \[\m(\tau) = m_1(\tau_1) \cdots m_{n-d}(\tau_{n-d}).\]
    Define the elements $\psi_{\m} \in \Hom_{\S_n}\left(\bigwedge^{n-d} \C^n, R_n^{(k,\ell)}\right)$ on the basis as follows:
    \[
    \psi_{\m}(e_\tau) = \cY_\tau(\m(\tau)). 
    \]
\end{definition}

\begin{proposition}\label{prop:psi-m-well-defined}
    The elements $\psi_{\m}$ are $\S_n$-equivariant and thus well defined.
 Moreover, for any subset $\tau \subset [n]$ of size $n-d$, we have that, 
 \[ \psi_{\m}(e_\tau)  =  \frac{1}{(n-d)!} \sum_{\sigma \in \S_{n-d}} \sgn(\sigma ) m_1( \tau_{\sigma(1)}) \cdots m_{n-d}(\tau_{\sigma(n-d)}).\]
\end{proposition}
\begin{proof}
Recall that for any $\sigma \in \S_n$, we have that $\sigma \cY_\tau = \cY_{\sigma(\tau)} \sigma$. Thus, 
\[
\sigma \psi_{\m}(e_\tau) = \sigma \cY_\tau (\m(\tau)) = \cY_{\sigma(\tau)} \sigma(\m(\tau)) = \cY_{\sigma(\tau)} \m(\sigma(\tau)) = \psi_{\m}(e_{\sigma(\tau)}) =  \psi_{\m}(\sigma e_\tau).
\]    

So then for a fixed ordered subset $\tau$ let $w\in \S_n$ be any permutation such that $w(\tau_i)=i$, so that $w(\tau)=[n-d]$ and $\cY_\tau = w^{-1} \cY_{ w(\tau) } w =  w^{-1} \cY_{[ n-d] } w$. 
Then, 
\begin{align*}
    \psi_{\m}(e_\tau) 
    &= \cY_\tau(\m(\tau)) 
    = w^{-1} \cY_{[n-d]} w(\m(\tau))\\
    & = \frac{1}{(n-d)!} \sum_{\sigma \in \S_{n-d}} 
    \sgn(\sigma ) m_1(w^{-1}\sigma w( \tau_1)) \cdots m_{n-d}(w^{-1} \sigma w(\tau_{n-d})) \\
    &= \frac{1}{(n-d)!} \sum_{\sigma \in \S_{n-d}} 
    \sgn(\sigma ) m_1( \tau_{\sigma(1)}) \cdots m_{n-d}(\tau_{\sigma(n-d)}),
\end{align*}
as claimed.
\end{proof}

We now prove that the monomial antisymmetrization operators 
$\psi_{\m}$ span $\Hom_{\S_n}\left(\bigwedge^{n-d}\C^n, R_n^{(k, \ell)}\right)$. 

\begin{theorem}\label{thm:spanning} 
    Suppose $0\leq d \leq n-1$ and 
    fix $\beta \subseteq [n]$ with $|\beta|=n-d$. 
    The subspace 
    $\cY_{\beta} \cC_{\overline{\beta}}\left(R_n^{(k, \ell)}\right)$ 
    is spanned by the set 
    \[
        \left\{\psi_{\m}(e_\beta) \mid \m = m_1 \cdots m_{n-d},  
        \text{ for some monomials } 
        m_i(j) \in R_n^{(k, \ell)}\right\}.
    \]
    Thus, $\Hom_{\S_n}\left(\bigwedge^{n-d}\C^n, R_n^{(k, \ell)}\right)$ 
    is spanned by 
    \[
        \left\{\psi_{\m} \mid \m = m_1 \cdots m_{n-d}, 
        \text{ for some monomials } 
        m_i(j) \in R_n^{(k, \ell)}\right\}.
    \]
\end{theorem}

\begin{proof}
    For this proof, when $d=0$, extend the necessary definitions 
    as needed to allow for $\overline{\beta} = \varnothing$.

    Recall from equation~\eqref{eq:quotient-map} the quotient 
    map $\Phi: \C_n^{(k,\ell)} \to R_n^{(k,\ell)}$. 
    By $\S_n$-equivariance, $\Phi$ commutes with $\cY_\beta\cC_{\overline{\beta}}$.
    Hence, for any $r \in \cY_{\beta} \cC_{\overline{\beta}}\left(R_n^{(k,\ell)}\right)$ 
    and any $r' \in \C_n^{(k,\ell)}$ satisfying $\Phi(r') = r$, 
    define $\tilde{r} \coloneqq \cY_{\beta} \cC_{\overline{\beta}}(r')$. 
    Then $\tilde{r}$ satisfies $\Phi(\tilde{r}) = 
    \cY_{\beta} \cC_{\overline{\beta}}(\Phi(r'))= \cY_{\beta} \cC_{\overline{\beta}}(r) = r $, 
    with $(\sigma,\tau)\tilde{r} = \sgn(\sigma)\tilde{r}$ for all 
    $(\sigma,\tau) \in \S_\beta \times \S_{\overline{\beta}}$.

    So then, for $r$ and $\tilde{r}$ as above, we factor $\tilde{r}$ 
    into products of polynomials in disjoint variables as follows.
    Write $\tilde{r} = \sum_a v_a(\beta) u_a(\overline{\beta})$ where 
    $v_a(\beta) = v_a\left(\bm{x}_\beta^{(\bullet)}, \bm{\theta}_\beta^{(\star)}\right)$ 
    is a monomial in variables $x_i^{(j)}$, $\theta_i^{(j')}$ 
    with $i \in \beta$ and $j \in [k], j' \in [\ell]$, and where 
    $u_a(\overline{\beta})=u_a\left(\bm{x}_{\overline{\beta}}^{(\bullet)}, \bm{\theta}_{\overline{\beta}}^{(\star)}\right)$ is a polynomial in variables $x_i^{(j)}$, $\theta_i^{(j')}$ with $i \in {\overline{\beta}}$ and $j \in [k], j' \in [\ell]$. In particular, since $\tilde{r} = \cC_{\overline{\beta}} \tilde{r} = \sum_a v_a(\beta)\cC_{\overline{\beta}} ( u_a(\overline{\beta}))$, then it follows that $\cC_{\overline{\beta}} ( u_a(\overline{\beta}))= u_a(\overline{\beta})$ so that $u_a(\overline{\beta}) \in \left(\C_{\overline{\beta}}^{(k,\ell)}\right)^{\S_{\overline{\beta}}}$. 

    We now rewrite the product $v_a(\beta) u_a(\overline{\beta})$ 
    using only the variables indexed by $\beta$. 
    Let $u_a(\overline{\beta}) = u + u_a(\overline{\beta})'$ with 
    $u \in \C$ constant and 
    $u_a(\overline{\beta})' \in 
    \left(\C_{\overline{\beta}}^{(k,\ell)}\right)_+^{\S_{\overline{\beta}}}$ 
    having constant term zero, then 
    $v_a(\beta) u_a(\overline{\beta}) = v_a(\beta) (u + u_a(\overline{\beta})') = 
    v_a(\beta) u + v_a(\beta) u_a(\overline{\beta})'$ where $v_a(\beta) u$ is 
    already written with only variables indexed by $\beta$. 
    Thus, it suffices to consider $u_a(\overline{\beta}) \in 
    \left(\C_{\overline{\beta}}^{(k,\ell)}\right)_+^{\S_{\overline{\beta}}}$, 
    which by Proposition ~\ref{prop:ideal-generators} 
    is a polynomial in the polarized power sums.

    Now, since $\beta$ and $\overline{\beta}$ partition $[n]$, 
    and $p_{\r;\s}([n]) \in I_n^{(k,\ell)}$ for $(\r,\s) \neq (\underline{0},\underline{0})$, 
    \[
        p_{\r;\s}(\overline{\beta}) = p_{\r;\s}([n]) - p_{\r;\s}({\beta})
        \equiv -p_{\r;\s}({\beta})\pmod {I_n^{(k,\ell)}}.
    \]
    Thus, $u_a(\overline{\beta}) \equiv u'_a(\beta) \pmod {I_n^{(k,\ell)}}$ 
    for some polynomial $u'_a(\beta)$.

    Thus, passing to the quotient, since $r$ is $\S_\beta$-antisymmetric 
    it follows that in $R_n^{(k,\ell)}$:
    \[
        r = \cY_\beta(r) 
        = \cY_\beta \left( \sum_{a} v_a(\beta) u'_a(\beta)\right) 
        = \sum_a \cY_\beta(v_a(\beta) u'_a(\beta)).
    \]
    Now, since both $v_a(\beta)$ and $u'_a(\beta)$ are polynomials, 
    then for each $a$ there exist monomials 
    \[
        \m_{a,b} = (m_{a,b})_1 \cdots (m_{a,b})_{n-d}
    \]
    such that the product $v_a(\beta) u'_a(\beta) =\sum_b c_{a,b} \m_{a,b}(\beta)$ 
    for some coefficients $c_{a,b} \in \C$. Hence,
    \[
        r = \sum_a \cY_\beta(v_a(\beta) u'_a(\beta)) 
        = \sum_a \sum_b c_{a,b} \cY_\beta( \m_{a,b}(\beta))
        = \sum_a \sum_b c_{a,b} \psi_{\m_{a,b}}(e_\beta).  
    \]
    Thus, $r$ is in the span of $\{\psi_{\m}(e_\beta)\}$ as desired. 
    The final claim now follows by Proposition~\ref{prop:hook-hom-iso}.
\end{proof}

Finally, as claimed above, we prove that the monomial antisymmetrization operators 
$\psi_{\m} \in \Hom_{\S_n}\left(\bigwedge^{n-d} \C^{n}, R_n^{(k,\ell)}\right)$ vanish 
when applied to any nonconstant monomial in $R_n^{(k,\ell)}$. 
As we shall see, the following Lemma is key in proving that 
$\Psi_{f \cdot p_{\r;\s;1}} = 0$ and thus the kernel of $\pi^{\sgn}$ is trivial.

\begin{lemma}\label{lem:Zero-sum}  Let $0\leq d \leq n-1$. 
    For any $\psi \in \Hom_{\S_n}\left(\bigwedge^{n-d}\C^n, R_n^{(k,\ell)}\right)$ and fixed monomial 
    $m_{\r;\s}(i) \in R_n^{(k,\ell)}$ with $(\r,\s) \neq (\underline{0},\underline{0})$, 
    we have that for all $\zeta \in \bigwedge^{n-d-1}\C^n$,
    \[
        \sum_{i=1}^n \psi(e_i \wedge \zeta) m_{\r;\s}(i) = 0. 
    \]
\end{lemma}

\begin{proof}
    By linearity, it suffices to check on the basis, so we may assume $\zeta = e_\mu$ for a subset $\mu = \{\mu_1 < \cdots < \mu_{n-d-1}\} \subseteq [n]$. 
    Likewise, from Theorem \ref{thm:spanning}, by linearity it suffices to consider the case $\psi = \psi_{\m}$ for some product monomial $\m=m_1\cdots m_{n-d}$. So then, fix $i \in [n] \setminus \mu$ and set $(\beta_1,\ldots, \beta_{n-d}) = (i,\mu_1,\ldots,\mu_{n-d-1})$. By Proposition~\ref{prop:psi-m-well-defined},
    \begin{align*}
        \psi_{\m}(e_i \wedge e_\mu) 
        &= \frac{1}{(n-d)!} \sum_{\sigma \in \S_{n-d}} \sgn(\sigma) m_1(\beta_{\sigma(1)}) \cdots m_{n-d}(\beta_{\sigma({n-d})}) \\
        &= \frac{1}{(n-d)!}\sum_{a=1}^{n-d} \varepsilon_a \left(\sum_{\substack{\sigma \in \S_{n-d}\\ \sigma(a)=1}} \sgn(\sigma) m_1(\beta_{\sigma(1)}) \cdots \widehat{m_a(\beta_1)}\cdots m_{n-d}(\beta_{\sigma({n-d})})\right) m_a(\beta_1)
    \end{align*}
    where $\varepsilon_a = \pm 1$ is the Koszul sign, and $\widehat{m_a(\beta_1)}$ denotes that this term has been removed. Thus, for each $1\leq a \leq n-d$ set
    \[
    B_a\coloneqq \frac{\varepsilon_a}{(n-d)!} \sum_{\substack{\sigma \in \S_{n-d}\\ \sigma(a)=1}} \sgn(\sigma) m_1(\beta_{\sigma(1)}) \cdots \widehat{m_a(\beta_1)}\cdots m_{n-d}(\beta_{\sigma({n-d})})
    \]
    so that $\psi_{\m}(e_i \wedge e_\mu) = \sum_{a=1}^{n-d} B_a m_a(i)$ since $\beta_1=i$. 

    Note that $(n-d)!\sum_{a=1}^{n-d} B_a m_a(i) = 
    \sum_{\sigma \in \S_{n-d}} \sgn(\sigma) m_1(\beta_{\sigma(1)}) \cdots 
    m_{n-d}(\beta_{\sigma({n-d})})$ is defined for every $i \in [n]$, 
    not just $i \in [n]\setminus \mu$. 
    Indeed, suppose $i \in \mu$ so that $\beta_a = \beta_1=i$ for some $a \neq 1$.
    Because for any $\sigma \in \S_{n-d}$ with $\sigma(b) = a$ and $\sigma(c)=1$ there exists 
    $\nu=(1\,a) \sigma \in \S_{n-d}$ satisfying 
    $\beta_{\sigma(j)} = \beta_{\nu(j)}$ for all $j \neq b,c$, with 
    $\beta_{\sigma(b)}=\beta_{a}=\beta_1=\beta_{\nu(b)}$ and 
    $\beta_{\sigma(c)} = \beta_1=\beta_a = \beta_{\nu(c)}$, then 
    \[
        \prod_{j=1}^{n-d} m_j(\beta_{\sigma(j)}) = \prod_{j=1}^{n-d} m_j(\beta_{\nu(j)}). 
    \]
    Since $\sgn(\nu)=-\sgn(\sigma)$ then 
    \[
        (n-d)! \sum_{a=1}^{n-d} B_a m_a(i) = 
        \sum_{\sigma \in \S_{n-d}} 
        \sgn(\sigma) m_1(\beta_{\sigma(1)}) \cdots m_{n-d}(\beta_{\sigma({n-d})})=0.
    \]
    On the other hand, because $e_i \wedge e_\mu = 0$ for all $i \in \mu$, 
    then $\psi_{\m}(e_i \wedge e_\mu)$ vanishes, so the identity 
    $\psi_{\m}(e_i \wedge e_\mu) = \sum_{a} B_a m_a(i)$ holds for all $i \in [n]$.
    Thus, summing over all $i$ we obtain, 
    \begin{align} \label{eq:Ba-sum}
        \sum_{i=1}^n \psi_{\m}(e_i \wedge e_\mu) m_{\r;\s}(i) 
        =  \sum_{i=1}^n \sum_{a=1}^{n-d} B_a m_a(i) m_{\r;\s}(i)
        = \sum_{a=1}^{n-d} B_a \sum_{i=1}^n  m_a(i) m_{\r;\s}(i).
    \end{align}
    
    Now, since the product of any two monomials $m_{\r';\s'}(i)m_{\r;\s}(i) = \pm m_{\r'+\r;\s'+\s}(i)$ with the sign independent of $i$, then for $m_a(i)= m_{\r^{(a)};\s^{(a)}}$, it follows that
    $\sum_{i=1}^n  m_a(i) m_{\r;\s}(i)=\pm p_{\r^{(a)}+\r; \s^{(a)}+\s}$, a polarized power sum, or zero. In particular, it is invariant under the action of $\S_n$. 
    Moreover, since $(\r,\s)\neq (\underline{0},\underline{0})$ then $1 \leq \deg(m_{\r;\s}(i))$ and thus $m_{\r;\s}(i)$ is nonconstant. Hence $\sum_{i=1}^n  m_a(i) m_{\r;\s}(i)$, and thus the sum in \eqref{eq:Ba-sum}, must lie inside $I_n^{(k,\ell)}$ so that 
    \[\sum_{i=1}^n \psi_{\m}(e_i \wedge e_\mu) m_{\r;\s}(i)=0\]
    in $R_n^{(k,\ell)}$ as desired.
\end{proof}

With all the ingredients in hand, we are now ready to state and prove the main result of this section. This generalizes the $(k,\ell) = (2,0)$ case from \cite{GLM}.

\begin{theorem}\label{thm:iso-sign}
    For any $n\geq 1$ and $k, \ell \geq 0$, the multiplication map $\pi: R_n^{(k,\ell)} \otimes R_n^{(0,1)} \to R_n^{(k,\ell+1)}$ is an isomorphism on sign-isotypic components. 
\end{theorem}

\begin{proof}
    By Corollary~\ref{cor:kernel-isotypic} and equation~\eqref{eq:sign-hook} we know that the kernel of $\pi^{\sgn}$ is graded  by $R_n^{(0,1)}$, so that $\ker(\pi)^{\sgn} = \bigoplus_{d=0}^{n-1} \ker(\pi)^{\sgn}_d$. 
    In degree $d=0$, since $\bigwedge\nolimits^0 \C^{n-1} \cong \C$, then $(\pi)_0 = \iota$ the inclusion map from equation \eqref{eq:ring-inclusion}. Pointedly, since $\iota$ is injective, then $(\pi)_0$ is an isomorphism. 

    So assume $d>0$ so that each graded piece $\ker(\pi)_d^{\sgn}$ is spanned by 
     \[
        \left\{  f \cdot p_{\r;\s;1}  \mid  \r \in (\mathbb{Z}_{\geq 0})^{k},\, 
    \s \in \{0,1\}^{\ell} \text{ with  } (\r,\s) \neq (\underline{0}, \underline{0}),\,
        f \in  \left(R_n^{(k,\ell)} \otimes \bigwedge\nolimits^{d-1} \C^{n-1}\right)^{\sgn} \right\}.
    \] 
   
    Now, recall that $\Hom_{\S_n}\left( \sgn, R_n^{(k,\ell)}  \otimes \bigwedge\nolimits^{d} \C^{n-1}\right) \cong \left(R_n^{(k,\ell)}  \otimes \bigwedge\nolimits^{d} \C^{n-1}\right)^{\sgn}$ via the identification $F \mapsto F(1)$, so that by Proposition \ref{prop:Hom-iso} there is an isomorphism
     \[
      \left(R_n^{(k,\ell)}  \otimes \bigwedge\nolimits^{d} \C^{n-1}\right)^{\sgn}
       \to \Hom_{\S_n}\left( \bigwedge\nolimits^{n-d-1}\C^{n-1},R_n^{(k,\ell)} \right)
    \]
    sending $g \mapsto \Psi_g$. Thus, for any 
    $ f \in  \left(R_n^{(k,\ell)} \otimes \bigwedge\nolimits^{d-1} \C^{n-1}\right)^{\sgn}$
    we have 
    \[
        f \cdot p_{\r;\s;1} =0 \text{ if and only if } \Psi_{f  \cdot p_{\r;\s;1}}=0.
    \]
    By linearity, it suffices to check $\Psi_{f \cdot p_{\r;\s;1}}=0$ 
    on every basis element $e_{\overline{\mu}} \in \bigwedge^{n-d-1}\C^{n-1}$. 
    So suppose that $f = \sum_{|\tau| = d-1} r_\tau \otimes e_\tau$ and $p_{\r;\s;1} = \sum_{i=1}^n m_{\r;\s}(i) \otimes e_i$.
    By the Koszul sign rule, 
    \[
        f \cdot p_{\r;\s;1} 
        = (-1)^{(d-1)|\s|}\sum_{|\tau|=d-1} \sum_{i=1}^n  r_\tau m_{\r;\s}(i) \otimes (e_\tau \wedge e_i).
    \]
    Then, by definition of $\Psi$, we write
    \begin{align*}
        (-1)^{(d-1)|\s|} \Psi_{f \cdot p_{\r;\s;1}}(e_{\overline{\mu}}) 
        =  \sum_{|\tau|=d-1} \sum_{i=1}^n  
        \langle e_\tau \wedge e_i , e_{\overline{\mu}} \rangle r_\tau m_{\r;\s}(i)
        =  \sum_{i=1}^n \left( \sum_{|\tau|=d-1} 
        \langle e_\tau \wedge e_i , e_{\overline{\mu}} \rangle r_\tau \right) m_{\r;\s}(i).
    \end{align*}
    Since $\langle e_\tau \wedge e_i,e_{\overline{\mu}} \rangle \omega 
    = (e_\tau \wedge e_i) \wedge e_{\overline{\mu}} 
    = e_\tau \wedge (e_i \wedge e_{\overline{\mu}}) 
    = \langle  e_\tau , e_i \wedge e_{\overline{\mu}} \rangle \omega$, 
    then
    \begin{align*}
        (-1)^{(d-1)|\s|} \Psi_{f \cdot p_{\r;\s;1}}(e_{\overline{\mu}}) 
        =\sum_{i=1}^n \left( \sum_{|\tau|=d-1}
        \langle  e_\tau , e_i \wedge e_{\overline{\mu}} \rangle
        r_\tau \right) m_{\r;\s}(i)
        = \sum_{i=1}^n \Psi_{f}(e_i \wedge e_{\overline{\mu}})  m_{\r;\s}(i),
    \end{align*}
    where $\Psi_f \in \Hom_{\S_n}\left(\bigwedge^{n-d}\C^{n-1}, R_n^{(k,\ell)}\right)$.

    So consider the $\S_n$ equivariant surjection 
    $\Omega: \bigwedge^{n-d}\C^n \twoheadrightarrow \bigwedge^{n-d}\C^{n-1}$. 
    Since $\Omega(e_i \wedge e_{\overline{\mu}})=  e_i \wedge e_{\overline{\mu}}$ 
    for all $i$ and $\overline{\mu}$, and by assumption 
    $(\r,\s) \neq (\underline{0}, \underline{0})$, then by Lemma \ref{lem:Zero-sum} 
    \[
        \sum_{i=1}^n \Psi_{f}(e_i \wedge e_{\overline{\mu}})  m_{\r;\s}(i) = 
        \sum_{i=1}^n \Psi_{f}(\Omega(e_i \wedge e_{\overline{\mu}}) ) m_{\r;\s}(i) = 0.
    \]
    Thus $\Psi_{f \cdot p_{\r;\s;1}} = 0$, hence $f \cdot p_{\r;\s;1} = 0$, 
    so by Corollary \ref{cor:kernel-isotypic} we have that $\pi^{\sgn}$ is an isomorphism.
\end{proof}

As an immediate consequence, we obtain a description for the sign character of 
$R_n^{(k,\ell+1)}$ in terms of the hook characters of $R_n^{(k,\ell)}$.

\begin{corollary}\label{cor:hook-to-sign-sf}
    For any $n\geq 1$ and $k,\ell \geq 0$, with $a$ recording the degree in the alphabet $\bm{\theta}^{(\ell+1)}$, we have
    \begin{align*}
        \left\langle \Frob(R_n^{(k,\ell+1)}; \mathbf{q};\mathbf{u}, a), s_{(1^n)} \right\rangle
   &= \left\langle \Frob(R_n^{(k,\ell)} \otimes R_n^{(0,1)}; \mathbf{q};\mathbf{u}, a), s_{(1^n)} \right\rangle\\
   &= \sum_{d=0}^{n-1} a^d \left\langle \Frob(R_n^{(k,\ell)}; \mathbf{q};\mathbf{u}), s_{(d+1,1^{n-d-1})}\right\rangle.
    \end{align*}
\end{corollary}

\begin{proof}
    The first equality follows directly from Theorem \ref{thm:iso-sign} by applying Frobenius characters and the second is equation~\eqref{eq:sign-hook} (see \cite{GLM}*{Prop. 3.6}).
\end{proof}

\subsection{The socle contains the sign component}
Recall that since $R_n^{(k,\ell)}$ is graded by total homogeneous degree, 
then $R_n^{(k,\ell)}$ is a graded local ring with maximal ideal 
\[
    (R_n^{(k,\ell)})_+= 
    \left\langle  x_i^{(\bullet)}, \theta_i^{(\star)} \mid  
    i \in [n], \bullet \in [k], \star \in [\ell] \right\rangle.
\]

Thus, the \newword{socle} of $R_n^{(k,\ell)}$ can be defined as 
the annihilator of $(R_n^{(k,\ell)})_+$ or, equivalently, as the set 
\[
    \soc(R_n^{(k,\ell)}) = \left\{ f \in R_n^{(k,\ell)} \mid f x_i^{(\bullet)} =0 
    \text{ and } 
    f \theta_i^{(\star)} =0 \text{ for all } i \in [n], \bullet \in [k], \star \in [\ell] \right\}.
\]

One of Haiman's most celebrated results is the proof of the $(n+1)^{n-1}$ theorem, which computes the dimension of the space of diagonal coinvariants $R_n^{(2,0)}$ \cite{Haiman2002}. In the same work he also proved his operator conjecture \cite{Haiman1994}*{Conj.~5.1.1}, now the \emph{operator theorem} \cite{Haiman2002}*{Thm.~4.2}, a standard tool in the study of diagonal coinvariants. One particularly remarkable consequence of the operator theorem, conjectured in \cite{Haiman1994}*{Conj.~5.4.1}, is that the socle of $R_n^{(2,0)}$ coincides with its sign isotypic component (see \cite{Haiman2002}*{Sec.~4}). 

In the theorem below we extend one direction of Haiman's result to all rings $R_n^{(k,\ell)}$. 
Namely, as a special case of Lemma \ref{lem:Zero-sum} we obtain that 
the alternating component of $R_n^{(k,\ell)}$ is always contained in 
the socle of $R_n^{(k,\ell)}$.

\begin{theorem}\label{thm:socle} 
    For all $n\geq 1$ and $k,\ell \geq 0$ we have that
    \[
        \left(R_n^{(k,\ell)}\right)^{\sgn} \subseteq \soc\left(R_n^{(k,\ell)}\right).
    \]
\end{theorem}

\begin{proof}
Denote by $\underline{\delta}_b = (0,\ldots, 0, 1, 0, \ldots, 0)$ the sequence with $1$ in the $b^{th}$ position and zeroes elsewhere. Then, it follows from the definition of $m_{\r;\s}(i)$ that,
\[
m_{\underline{\delta}_b;\underline{0}}(i) = x_i^{(b)} 
\qquad \text{ and }\qquad
m_{\underline{0};\underline{\delta}_b}(i) = \theta_i^{(b)}.
\]
From Lemma \ref{lem:Zero-sum} we know that for any $\psi \in \Hom_{\S_n}\left(\bigwedge^{n-d}\C^n, R_n^{(k,\ell)}\right)$ and any nonconstant monomial $m_{\r;\s}(i)$ we have, 
$
    \sum_{i=1}^n \psi(e_i \wedge \zeta) m_{\r;\s}(i) = 0 
$
for any $\zeta \in \bigwedge^{n-d-1}\C^n$. So then, take $d=0$ so that $\psi \in \Hom_{\S_n}\left(\sgn, R_n^{(k,\ell)}\right)$ and for each $1\leq a \leq n$ set $\zeta_a = e_1\wedge \cdots \wedge \widehat{e_a} \wedge \cdots \wedge e_n$ to be the wedge of all basis elements with $e_a$ omitted. Let $\omega_n = e_1 \wedge \cdots \wedge e_n \in \bigwedge \C^n$. Then by Lemma \ref{lem:Zero-sum}, for all $1\leq a \leq n$  and all $b \in [k]$ or $b \in [\ell]$, respectively, we have that, up to a sign,
\begin{align*}
\pm\psi(\omega_n) x_a^{(b)}  &= \sum_{i=1}^n \psi(e_i \wedge \zeta_a) x_i^{(b)}= \sum_{i=1}^n \psi(e_i \wedge \zeta_a) m_{\underline{\delta}_b;\underline{0}}(i) = 0, \\
\pm\psi(\omega_n) \theta_a^{(b)}  &= \sum_{i=1}^n \psi(e_i \wedge \zeta_a) \theta_i^{(b)}= \sum_{i=1}^n \psi(e_i \wedge \zeta_a) m_{\underline{0};\underline{\delta}_b}(i) = 0.
\end{align*}
Thus $\psi(\omega_n) \in \soc(R_n^{(k,\ell)})$ for all $\psi \in \Hom_{\S_n}\left(\sgn, R_n^{(k,\ell)}\right)$. 
The claim follows since \[\left(R_n^{(k,\ell)}\right)^{\sgn} = \left\{ \psi(\omega_n) \mid \psi \in \Hom_{\S_n}\left(\sgn, R_n^{(k,\ell)}\right)\right\}.\] 
\end{proof}

%%%%%%%%%%%%%%%%%%%%%%%%%%%%%%%%%%%%%%%%%%%%%%%%%%%%%%%%%%%%%%%%%%%%%%
\section{Sign components of \texorpdfstring{$R_n^{(k,\ell)} \otimes R_n^{(a,b)}$}{Rn(k,l)xRn(a,b)}}\label{sec:miniswaddle}
%%%%%%%%%%%%%%%%%%%%%%%%%%%%%%%%%%%%%%%%%%%%%%%%%%%%%%%%%%%%%%%%%%%%%%

From Theorem \ref{thm:iso-sign} one may ask if the result holds in greater generality; namely, are there other values of $(a,b)$ for which the sign-isotypic components of $R_n^{(k,\ell)} \otimes R_n^{(a,b)}$ and $R_n^{(k+a,\ell+b)}$ are isomorphic? Evidently, if $(a,b)=(0,0)$, the map is the identity and the question is moot. For nontrivial $(a,b)$, the answer depends on the value of $n$. 

When $n$ is very small (i.e., $n=1,2$) the rings degenerate and one gets an isomorphism on sign components for all values of $(a,b)$. 

The situation is far more interesting when $n\geq 3$, where we show that $\pi^{\sgn}$ is injective if and only if at least one of $(k,\ell)$ and $(a,b)$ is $(0,0)$ or $(0,1)$ (Theorem~\ref{thm:neg-result}). 
This implies that Theorem~\ref{thm:iso-sign} covers the only nontrivial case.

\subsection{The \texorpdfstring{$n\geq 3$}{n>=3} case}\label{subsec:greater-than-3} 
Our strategy to show that for $n\geq 3$ and $k,\ell$ arbitrary, the projection $\pi^{\sgn}: \left(R_n^{(k,\ell)}\otimes R_n^{(a,b)}\right)^{\sgn} \to \left(R_n^{(k+a,\ell+b)}\right)^{\sgn}$ fails to be an isomorphism for $(a,b) \not\in \{(0,1), (0,0)\}$, consists of dividing the argument into three distinct subcases where,
\begin{enumerate}
    \item  both $ (k,\ell), (a,b) \geq (1,0)$, 
    \item  $ (k,\ell)\geq (1,0)$ and $(a,b) \geq (0,2)$, or 
    \item both $ (k,\ell), (a,b) \geq (0,2)$.
\end{enumerate} 
By symmetry, the remaining situation when $ (k,\ell)\geq (0,2)$ and $(a,b) \geq (1,0)$ need not be checked. Since all three subcases rely on dimensionality arguments, we recall some useful facts about computing sign components. 

Recall that for any $\S_n$-modules $M \cong \bigoplus_{\lambda \vdash n} S^\lambda \otimes M_\lambda$ and $N \cong \bigoplus_{\mu \vdash n} S^\mu \otimes N_\mu$, with $M_\lambda$ and $N_\mu$ the respective multiplicity spaces, 
by tensor-hom adjunction we have that,
\begin{equation}
    \label{eq:tensor-sign decomposition}
    \begin{aligned} 
        \Hom_{\S_n}\left( \sgn, M \otimes N\right) 
        &\cong \Hom_{\S_n}\left( \sgn \otimes N, M \right) 
        \cong \bigoplus_{\substack{\lambda \vdash n \\ \mu \vdash n}} 
        \Hom_{\S_n}\left( \sgn \otimes S^\mu, S^\lambda \right)  \otimes M_\lambda \otimes N_\mu \\ 
        &\cong \bigoplus_{\substack{\lambda \vdash n \\ \mu \vdash n}} 
        \Hom_{\S_n}\left( S^{\mu^t}, S^\lambda \right) \otimes M_\lambda \otimes N_\mu 
        \cong \bigoplus_{\lambda \vdash n} M_\lambda \otimes N_{\lambda^t}.
    \end{aligned}
\end{equation}

In particular, if we consider $M = R_n^{(1,0)}$, 
since by \cite{Chevalley}, 
$R_n^{(1,0)} \cong \C[\S_n] \cong \bigoplus_{\lambda \vdash n} S^{\lambda}\otimes \C^{f^\lambda}$ 
is the regular representation of $\S_n$, 
where $f^\lambda$ is the number of standard Young tableaux of shape $\lambda$, 
and since (as vector spaces) $S^\lambda \cong \C^{f^\lambda}$ with $f^\lambda = f^{\lambda^t}$, 
then by equation~\eqref{eq:tensor-sign decomposition}:
\begin{equation}
    \label{eq:tensor-with-(1,0)}
    \Hom_{\S_n}\left( \sgn, R_n^{(1,0)} \otimes N\right) 
    \cong 
    \bigoplus_{\lambda \vdash n} \C^{f^\lambda} \otimes N_{\lambda^t}
    \cong \bigoplus_{\lambda^t \vdash n} S^{\lambda^t} \otimes N_{\lambda^t} \cong N. 
\end{equation}

We begin by proving the first subcase. 

\begin{proposition} \label{prop:(1,0)x(1,0)}
    For any integers $n\geq 3$ and $a,k\geq 1$, $\ell,b\geq0$, the sign isotypic components of $R_n^{(k,\ell)} \otimes R_n^{(a,b)}$ and $R_n^{(k+a, \ell+b)}$ are not isomorphic. 
\end{proposition}

\begin{proof}
    By Lemma \ref{lem:tensor-commutative-diagram} for any $k,a \geq 1$ 
    it suffices to consider $(k,\ell)=(a,b)=(1,0)$.
    Then, by equation~\eqref{eq:tensor-with-(1,0)},
    \[
        \dim\left( \left(R_n^{(1,0)} \otimes R_n^{(1,0)}\right)^{\sgn}\right) 
        = \sum_{\lambda \vdash n} \left(f^\lambda \right)^2= n!.
    \]
    On the other hand, by \cite{Haiman2002}, 
    $\dim \left( \left(R_n^{(2,0)}\right)^{\sgn}\right) =\frac{1}{n+1}\binom{2n}{n}$, 
    the $n$th Catalan number. Since for any $n\geq 3$ we have that 
    $n!>\frac{1}{n+1}\binom{2n}{n}$, then the result follows.
\end{proof}

The second subcase follows in a similar fashion to the first.

\begin{proposition}
    \label{prop:(1,0)x(0,2)}
    For any integers $n\geq 3$ and $k\geq 1$, $b\geq 2$, $\ell,a\geq0$, the sign isotypic components of $R_n^{(k,\ell)} \otimes R_n^{(a,b)}$ and $R_n^{(k+a, \ell+b)}$ are not isomorphic. 
\end{proposition}

\begin{proof}
    By Lemma \ref{lem:tensor-commutative-diagram} it suffices to consider $(k,\ell)=(1,0)$ and $(a,b)=(0,2)$. 
    Then, by equation~\eqref{eq:tensor-with-(1,0)} and \cite{KimRhoades2022}*{Thm. 6.1} we have, 
    \[
        \dim\left( \left(R_n^{(1,0)} \otimes R_n^{(0,2)}\right)^{\sgn}\right) = \dim\left(R_n^{(0,2)}\right) = \binom{2n-1}{n}.
    \]
    As we shall see, Corollary~\ref{cor:12-sign} (which follows from Corollary~\ref{cor:hook-to-sign-sf} and \cite{BCJLR}*{Proposition 4.10}, so does not depend on this section) gives $\dim\left(\left(R_n^{(1,2)}\right)^{\sgn}\right) = 3^{n-1}$. Since for $n\geq 3$ we have $\binom{2n-1}{n}> 3^{n-1}$, we are done.
\end{proof}

The final case requires a bit more finesse, 
as the basic argument used above is no longer accessible to us. 
We begin recalling some useful combinatorial facts. 

For $n \in \Z_{\geq 0}$, the $(u,v)$-integers are the polynomials
\[[n]_{u,v} \coloneqq 
\frac{u^n-v^n}{u-v} = 
u^{n-1} + u^{n-2}v + \cdots + uv^{n-2} + v^{n-1} .\]
The second author computed the following characters of $R_n^{(0,2)}$, building on the work of Kim and Rhoades \cite{KimRhoades2022}. 
\begin{proposition}[{\cite{Lentfer2025}*{Thm. 19}}]
    \label{prop:02-all-characters}
    Let $\lambda = (\lambda_1, \lambda_2, 2^{c}, 1^{n-2c-\lambda_1-\lambda_2})$ with $\lambda_1 \geq \lambda_2 \geq 2$ and $c \geq 0$, and let $m \coloneqq n - 2c - \lambda_1 - \lambda_2 + 2$. If $\lambda_1 = \lambda_2$, then
    \begin{equation}\label{eq:(0,2)-fathook-character}
        \left\langle \Frob(R_n^{(0,2)};u,v), s_{\lambda} \right\rangle = (uv)^{c+\lambda_2-1}\big(uv[m-2]_{u,v} + [m-1]_{u,v} + [m]_{u,v}\big),
    \end{equation}
    and if $\lambda_1 > \lambda_2$, then
    \begin{equation}\label{eq:(0,2)-doublehook-character}
        \left\langle \Frob(R_n^{(0,2)};u,v), s_{\lambda} \right\rangle = (uv)^{c+\lambda_2-1}\big(uv[m-2]_{u,v} + (uv+1)[m-1]_{u,v} + [m]_{u,v}\big).
    \end{equation}
\end{proposition}

We can now prove the final subcase.

\begin{proposition}\label{prop:(0,2)x(0,2)}
    For any integers $n\geq 3$ and $\ell, b\geq 2$, $k,a\geq 0$, the sign isotypic components of $R_n^{(k,\ell)} \otimes R_n^{(a,b)}$ and $R_n^{(k+a, \ell+b)}$ are not isomorphic. 
\end{proposition}

\begin{proof}
By Lemma \ref{lem:tensor-commutative-diagram} it suffices to consider $(k,\ell)=(a,b)=(0,2)$. We consider two cases. 

First, suppose $n=3$. From direct computation one can show that 
\[\dim\left(\left(R_3^{(0,2)} \otimes R_3^{(0,2)}\right)^{\sgn}\right) = 15 > 14 = \dim\left((R_3^{(0,4)}\right)^{\sgn}).\]

Thus, assume $n \geq 4$ and let $\lambda = (n-2,2)$ so that $\lambda^t =(2,2,1^{n-4})$. 
By Proposition~\ref{prop:02-all-characters} 
applied to $\lambda$ (with $c = 0$ and $m = 2$; here $\lambda_1 = \lambda_2$ if $n = 4$ and $\lambda_1 > \lambda_2$ if $n \geq 5$), the presence of the term $(uv)[2]_{u,v} = uv(u+v)$ shows that $u^2v$ occurs in the bigraded Hilbert series of $(R_n^{(0,2)})^{(n-2,2)}$. 
Applied to $\lambda^t$ (with $c = 0$ and $m = n-2$), the presence of the term $uv[n-2]_{u,v}$ shows that $u^{n-2}v$ occurs in the bigraded Hilbert series of $(R_n^{(0,2)})^{(2,2,1^{n-4})}$.
Hence, the following homogeneous degree components are nonzero:
\[
\left(R_n^{(0,2)}\right)^{(n-2,2)}_{(2,1)} \neq 0 \qquad \text{and} \qquad \left(R_n^{(0,2)}\right)^{(2,2,1^{n-4})}_{(n-2,1)} \neq 0.
\]
Consequently, $\left(R_n^{(0,2)} \otimes R_n^{(0,2)}\right)^{\sgn}_{(2,1,n-2,1)} \neq 0$ since by equation~\eqref{eq:tensor-sign decomposition} we know that
\begin{align*}
        \dim \left(R_n^{(0,2)} \otimes R_n^{(0,2)}\right)^{\sgn}_{(2,1,n-2,1)} 
        \geq \dim \left(R_n^{(0,2)}\right)^{(n-2,2)}_{(2,1)} \cdot \dim \left(R_n^{(0,2)}\right)^{(2,2,1^{n-4})}_{(n-2,1)} >0.
\end{align*}

On the other hand, by \cite{KimRhoades2022}*{Thm. 6.1}, $R_n^{(0,2)}$ vanishes in all bidegrees $(i,j)$ with $i+j\geq n$. Since the map $\pi: R_n^{(0,2)} \otimes R_n^{(0,2)} \to R_n^{(0,4)}$ is a surjection and multidegree-preserving, then $R_n^{(0,4)}$ vanishes in all multidegrees $(d_1,d_2,d_3,d_4)$ where $d_1+d_2\geq n$ or $d_3+d_4 \geq n$. By Lemma \ref{lem:degree-permutation}, this implies $R_n^{(0,4)}$ vanishes in all multidegrees in which the sum of any two entries is at least $n$. Consequently, $(R_n^{(0,4)})_{(2,1,n-2,1)}=0$ and thus,
\[
\left(R_n^{(0,2)} \otimes R_n^{(0,2)}\right)^{\sgn}_{(2,1,n-2,1)}  \neq 0= \left(R_n^{(0,4)}\right)^{\sgn}_{(2,1,n-2,1)}.
\]
\end{proof}

By combining Theorem \ref{thm:iso-sign} with 
Propositions \ref{prop:(1,0)x(1,0)}, \ref{prop:(1,0)x(0,2)}, and \ref{prop:(0,2)x(0,2)}, 
the following theorem is now immediate.

\begin{theorem}\label{thm:neg-result}
    Let $n \geq 3$.
    The surjection
    \begin{equation*}
        \pi: R_{n}^{(k,\ell)} \otimes R_{n}^{(a,b)} \twoheadrightarrow R_{n}^{(k+a,\ell+b)},
    \end{equation*}
     is an isomorphism on sign-isotypic components if and only if at least one of $(k,\ell)$ or $(a,b)$ is $(0,0)$ or $(0,1)$.
\end{theorem}

\begin{proof}
    If $(a,b)$ or $(k,\ell)$ is $(0,0)$, then $\pi$ is the identity map. 
    If $(a,b)$ or $(k,\ell)$ is $(0,1)$, then Theorem~\ref{thm:iso-sign} asserts $\pi$ is an isomorphism on sign-isotypic components.
    The converse follows from Propositions \ref{prop:(1,0)x(1,0)}, \ref{prop:(1,0)x(0,2)}, and \ref{prop:(0,2)x(0,2)}.
\end{proof}

\subsection{The \texorpdfstring{$n=1$}{n=1} or \texorpdfstring{$n=2$}{n=2} case}\label{subsec:1-or-2}

When $n =1$ or $n=2$, the map $\pi$ is an isomorphism on sign-isotypic components 
for every choice of $(k,\ell)$ and $(a,b)$; 
thus the hypothesis $n \geq 3$ in Theorem~\ref{thm:neg-result} is necessary.

\begin{proposition}\label{prop:small-n-mini}
    Fix $n \in \{1,2\}$ and $k,\ell,a,b \geq 0$.
    Then the surjection $\pi: R_n^{(k,\ell)} \otimes R_n^{(a,b)} \twoheadrightarrow R_n^{(k+a,\ell+b)}$ is an isomorphism on sign-isotypic components.
\end{proposition}

\begin{proof}
    For $n=1$, every polynomial is $\S_1$-invariant, so for all $k,\ell,a,b$, $R_1^{(k,\ell)} = \C = R_1^{(a,b)}$, and thus $\pi:\C \otimes \C \twoheadrightarrow \C$ is an isomorphism.

    Let $n=2$. 
    Let $\bm{y} = {\{y_1,y_2\}}$ be any alphabet (bosonic or fermionic) of $\C_2^{(k,\ell)}$. 
    By Proposition~\ref{prop:ideal-generators}, $y_1 + y_2$ is in $I_2^{(k,\ell)}$, the defining ideal of $R_2^{(k,\ell)}$.
    One may verify that the set $\{1\} \cup \left\{y_1-y_2 : \bm{y} \text{ an alphabet in } 
    \C_2^{(k,\ell)}\right\}$ is a basis for $R_2^{(k,\ell)}$.
    Consequently,
    \begin{equation*}
        \Frob\left(R_2^{(k,\ell)}; \mathbf{q}; \mathbf{u}\right) = 
        s_{(2)} + \left(|\mathbf{q}| + |\mathbf{u}|\right)\, s_{(1,1)}.
    \end{equation*}

    Similarly, given $\bm{y}' = {\{y_1',y_2'\}}$ any alphabet (bosonic or fermionic) of 
    $\C_2^{(a,b)}$, then the set 
    $\{1\} \cup \left\{y_1'-y_2' : \bm{y}' \text{ an alphabet in } \C_2^{(a,b)}\right\}$ 
    is a basis for $R_2^{(a,b)}$. 
    Thus the sign-isotypic component of $R_2^{(k,\ell)} \otimes R_2^{(a,b)}$ has a basis consisting of elements $(y_1-y_2) \otimes 1$ and $1 \otimes (y_1'-y_2')$.  
    The map $\pi$ sends these to the elements $y_1-y_2$ and $y_1'-y_2'$; 
    since $\left\{y_1-y_2 : \bm{y} \text{ an alphabet in } \C_2^{(k,\ell)}\right\} \cup 
    \left\{y_1'-y_2' : \bm{y}' \text{ an alphabet in } \C_2^{(a,b)}\right\}$ 
    is a basis for the sign-isotypic component of $R_2^{(k+a,\ell+b)}$, 
    then $\pi$ is an isomorphism. 
\end{proof}

%%%%%%%%%%%%%%%%%%%%%%%%%%%%%%%%%%%%%%%%%%%%%%%%%%%%%%%%%%%%%%%%%%%%%%
\section{Other isotypic components of \texorpdfstring{$R_n^{(k,\ell)} \otimes R_n^{(0,1)}$}{Rn(k,l)xRn(0,1)}}\label{sec:mezzoswaddle}
%%%%%%%%%%%%%%%%%%%%%%%%%%%%%%%%%%%%%%%%%%%%%%%%%%%%%%%%%%%%%%%%%%%%%%

In this section, we classify the isotypic components on which the map $\pi: R_n^{(k,\ell)} \otimes R_n^{(0,1)} \twoheadrightarrow R_n^{(k,\ell+1)}$ fails to be injective.
The goal is to prove the following result.
\begin{theorem}\label{thm:mezzo-swaddle}
    Fix $n \geq 1$, $k \geq 1$, and $\ell \geq 0$, and let $\nu \vdash n$.
    Then the map
    \begin{equation*}
        \pi^\nu: \left(R_n^{(k,\ell)} \otimes R_n^{(0,1)}\right)^\nu \twoheadrightarrow \left(R_n^{(k,\ell+1)}\right)^\nu
    \end{equation*}
    is injective if and only if $\nu = (1^n)$.
\end{theorem}

Recall that $\bigwedge^d \C^{n-1}$ is the degree $d$ component of 
$R_n^{(0,1)}$ with basis $\{e_\tau\mid \tau \subseteq [n-1]\;,\; |\tau|=d\}$ and, 
as before, denote $\omega= e_1\wedge \cdots \wedge e_{n-1}$ so that 
$\bigwedge^{n-1}\C^{n-1} = \C \omega$ is the sign representation. 
For each $0\leq d\leq n-1$ set,
\[
    K_d\coloneqq \ker(\pi) \cap \left(R_n^{(k,\ell)}\otimes \bigwedge\nolimits^d \C^{n-1}\right).
\]

\begin{lemma}\label{lem:K_{n-1}}
    Fix $n \geq 1$ and $k,\ell \geq 0$. Let $\mathfrak{m} =(R_n^{(k,\ell)})_+$ 
    denote the maximal ideal in $R_n^{(k,\ell)}$.
    Then
    \begin{equation*}
        K_{n-1} = \mathfrak{m} \otimes \C \omega.
    \end{equation*}
    Hence, for all $\nu \vdash n$, the $\nu$-isotypic component of $K_{n-1}$ is 
    $\mathfrak{m}^{\nu^t} \otimes \C \omega$, where $\mathfrak{m}^{\nu^t}$ 
    denotes the $\nu^t$-isotypic component of $\mathfrak{m}$.
\end{lemma}

\begin{proof}
    Consider the map $\rho: R_n^{(k,\ell+1)} \twoheadrightarrow \bigwedge \C^{n-1}$ 
    sending $\theta_i^{(\ell+1)} \mapsto e_i$ and $x_i^{(a)}, \theta_i^{(b)} \mapsto 0$ 
    for all $1 \leq i\leq n$, $1\leq a \leq k$, and $1\leq b \leq \ell$. 
    Then for any $f \otimes g \in R_n^{(k,\ell)} \otimes \bigwedge \C^{n-1}$ we have
    \[
        \rho \circ \pi(f \otimes g) = \rho(fg) = \mathrm{CT}(f)g,
    \]
    where $\mathrm{CT}(f)$ denotes the constant term of $f$. Hence, 
    \[
        \ker(\pi) \subseteq \ker(\rho \circ \pi) = \mathfrak{m} \otimes \bigwedge \C^{n-1},
    \]
    and thus $K_{n-1} \subseteq \mathfrak{m} \otimes \C \omega$.

    Conversely, let $\bm{z}=\{z_1,\ldots, z_n\}$ with 
    $\bm{z} \in \{ \bm{x}^{(1)},\ldots, \bm{x}^{(k)}, 
    \bm{\theta}^{(1)}, \ldots, \bm{\theta}^{(\ell)}\}$. 
    Consider 
    \[
        F = \sum_{i=1}^n z_i \otimes e_i \in R_n^{(k,\ell)} \otimes \bigwedge \C^{n-1}
    \]
    where $e_n= -(e_1 + \cdots + e_{n-1})$. 
    Then, $\pi(F) =  \sum_{i=1}^n z_i\theta_i^{(\ell+1)}$ is clearly an 
    $\S_n$-invariant of positive degree and thus contained in $I_n^{(k,\ell+1)}$. 
    Hence, $F \in \ker(\pi)$.

    So then, for $\tau_a\coloneqq[n-1]\setminus \{a\}$ for $1 \leq a \leq n-1$, 
    set $e_{\overline{a}}\coloneqq e_{{\tau_{a}}} \in \bigwedge^{n-2}\C^{n-1}$, 
    so that for each $1\leq i \leq n-1$,
    \[
        e_i \wedge e_{\overline{a}} = \delta_{i,a} (-1)^{a-1} \omega
        \qquad \text{and} \qquad
        e_n \wedge e_{\overline{a}} = (-1)^a \omega.
    \]
    Since $F \in \ker(\pi)$ with $\ker(\pi)$ an ideal, then for all $1\leq a \leq n-1$,
    \[
        F ( 1 \otimes e_{\overline{a}}) = \sum_{i=1}^n z_i \otimes e_i\wedge e_{\overline{a}} 
        = (-1)^{a-1}(z_a-z_n) \otimes \omega \in \ker(\pi).
    \]
    Since $\sum_{i=1}^n z_i$ is $\S_n$-invariant, it must be contained in 
    $I_n^{(k,\ell)}$, so that $z_n = - \sum_{i=1}^{n-1} z_i$. Hence, 
    \[
        \sum_{a=1}^{n-1} (z_a -z_n)\otimes \omega = -n (z_n \otimes \omega) \in \ker(\pi)
    \]
    and so $z_a \otimes \omega = (z_a-z_n)\otimes \omega + z_n\otimes \omega \in \ker(\pi)$ 
    for all values of $a$. 
    Thus, $\langle \bm{z} \rangle \otimes \C\omega \subset \ker(\pi)$ for all 
    $\bm{z} \in \{ \bm{x}^{(1)},\ldots, \bm{x}^{(k)}, 
    \bm{\theta}^{(1)}, \ldots, \bm{\theta}^{(\ell)}\}$ with 
    $\langle \bm{z} \rangle = \langle z_1,\ldots, z_n\rangle$. 
    In particular, since $\langle \bm{z} \rangle \otimes \C\omega \subset 
    R_n^{(k,\ell)} \otimes \bigwedge^{n-1}\C^{n-1}$ then 
    $\langle \bm{z} \rangle \otimes \C\omega \subset K_{n-1}$ for all $\bm{z}$.

    Now, since by definition $\mathfrak{m} = \langle \bm{x}^{(1)},\ldots, 
    \bm{x}^{(k)}, \bm{\theta}^{(1)}, \ldots, \bm{\theta}^{(\ell)}\rangle$ then, 
    as desired,
    \[
        \mathfrak{m} \otimes \C\omega = \sum_{\bm{z}} \langle \bm{z} \rangle 
        \otimes \C\omega \subseteq K_{n-1}.
    \]
    The final claim now follows from the fact that $\C\omega \cong \sgn$ and 
    $S^{\lambda} \otimes \sgn \cong S^{\lambda^t}$ for all $\lambda \vdash n$ 
    so that
    $(K_{n-1})^\nu = (\mathfrak{m}\otimes \C\omega)^\nu = \mathfrak{m}^{\nu^t}\otimes \C\omega$.
\end{proof}

We can now prove Theorem~\ref{thm:mezzo-swaddle}.

\begin{proof}[Proof of Theorem~\ref{thm:mezzo-swaddle}]
    The kernel of the map induced by $\pi$ on $\nu$-isotypic components is $\ker(\pi)^\nu$.
    By Theorem~\ref{thm:iso-sign}, $\ker(\pi)^{(1^n)} = 0$.
    
    Now let $\nu \neq (1^n)$ and write $\lambda \coloneqq \nu^t$, so that $\lambda \neq (n)$.
    By Lemma~\ref{lem:K_{n-1}}, 
    $\ker(\pi)^\nu \supseteq (K_{n-1})^\nu = \mathfrak{m}^{\lambda} \otimes \C\omega$, 
    so it only remains to show that $\mathfrak{m}^{\lambda} \neq 0$.
    Since $k \geq 1$, setting all of the alphabets other than $\bm{x}^{(1)}$ 
    to zero induces a surjection $\rho: R_n^{(k,\ell)} \twoheadrightarrow R_n^{(1,0)}$ 
    of graded $\S_n$-modules, which maps $\mathfrak{m}$ onto the 
    maximal ideal $\left(R_n^{(1,0)}\right)_+$.
    Since $R_n^{(1,0)} \cong \bigoplus_{\lambda \vdash n} S^\lambda \otimes \C^{f^\lambda}$, 
    then every $S^\lambda$ appears with nonzero multiplicity in $R_n^{(1,0)}$. 
    Moreover, by \cite{stanley-invariants-finite-groups}*{Prop. 4.11} each 
    $\lambda$-isotypic component only appears in degrees $d$ with 
    $n(\lambda) \leq d \leq \binom{n}{2}-n(\lambda^t)$, where $n(\lambda) = \sum_i (i-1)\lambda_i$. 
    Since $n(\lambda)=0$ if and only if $\lambda = (n)$, 
    then $\left(R_n^{(1,0)}\right)^\lambda \subseteq \left(R_n^{(1,0)}\right)_+$. 
    Thus, $\rho(\mathfrak{m}^\lambda) =  (R_n^{(1,0)})^\lambda \neq 0 $, 
    hence $\mathfrak{m}^\lambda \neq 0$.
\end{proof}

\begin{remark}\label{rem:k0}
    The hypothesis $k \geq 1$ in Theorem~\ref{thm:mezzo-swaddle} is used in the proof to show that $\mathfrak{m}^\lambda \neq 0$ for all $\lambda \neq (n)$, because $R_n^{(k,\ell)}$ surjects onto the regular representation $R_n^{(1,0)}$.
    When $k=0$, not every irreducible representation of $\S_n$ needs to appear in $R_n^{(0,\ell)}$: for example, only hooks appear in $R_n^{(0,1)}$ (see equation~\eqref{eq:R(0,1)=hooks}) and only shapes contained in double-hooks appear in $R_n^{(0,2)}$ (see \cites{KimRhoades2022, Lentfer2025, Rosas}).

    We can see an example of more complicated behavior with the map $\pi: R_n^{(0,1)} \otimes R_n^{(0,1)} \twoheadrightarrow R_n^{(0,2)}$.
    First, $\pi^\nu$ can be injective because its domain is zero: this happens at $n=9$ with the partition $\nu = (3,3,3)$.
    Second, when $\pi^\nu$ is not injective, its kernel does not necessarily attain the top degree component $K_{n-1}$: this happens at $n=4$ with the partition $\nu = (2,2)$, where Lemma~\ref{lem:K_{n-1}} gives $(K_3)^\nu = 0$, since $\nu^t = (2,2)$ is not a hook.
    However, $S^{(2,2)}$ occurs in the bidegree $(2,2)$ component of 
    $R_4^{(0,1)} \otimes R_4^{(0,1)}$, which is $S^{(2,1,1)} \otimes S^{(2,1,1)}$. 
    But $R_4^{(0,2)}$ vanishes in bidegree $(2,2)$ by \cite{KimRhoades2022}*{Thm. 6.1}, so $\pi^{(2,2)}$ is not injective as its kernel contains a copy of $S^{(2,2)}$ in bidegree $(2,2)$. 
    This means that a result analogous to Theorem~\ref{thm:mezzo-swaddle} at $k=0$ requires further study of the exact $\S_n$-module structure of the modules $R_n^{(0,\ell)}$. 
\end{remark}

%%%%%%%%%%%%%%%%%%%%%%%%%%%%%%%%%%%%%%%%%%%%%%%%%%%%%%%%%%%%%%%%%%%%%%
\section{Sign components of \texorpdfstring{$R_{\mathfrak{B}_n}^{(k,\ell)} \otimes R_{\mathfrak{B}_n}^{(0,1)}$}{RBn(k,l)xRBn(0,1)}}\label{sec:type-B}
%%%%%%%%%%%%%%%%%%%%%%%%%%%%%%%%%%%%%%%%%%%%%%%%%%%%%%%%%%%%%%%%%%%%%%

In this section we extend Theorem~\ref{thm:iso-sign} to the determinant character of the hyperoctahedral group $\mathfrak{B}_n$. We begin by recalling some basic constructions but refer the reader to  \cite{Macdonald}*{Chapter I, Appendix B}, \cite{GeckPfeiffer}, and \cite{JiangLentfer}*{Section 2} for additional details.

\subsection{Representations of the hyperoctahedral group}\label{subsec:B-basics}
Let $\mathfrak{B}_n = \{+1,-1\} \wr \mathfrak{S}_n$ denote the hyperoctahedral group consisting of signed permutations. 
We write elements of $\mathfrak{B}_n$ as $g = (\epsilon, \sigma)$, where $\epsilon \in \{+1,-1\}^n$ and $\sigma \in \mathfrak{S}_n$.

A \newword{bipartition} $(\lambda, \mu)$ of $n$ is an ordered pair of two partitions whose sizes $|\lambda|$ and $|\mu|$ sum to $n$.
The bipartitions $(\lambda, \mu)$ of $n$ index the irreducible representations $V^{(\lambda, \mu)}$ of $\mathfrak{B}_n$.
The \newword{trivial representation} of $\mathfrak{B}_n$ is $V^{((n),\varnothing)}$.

Let $\Lambda_{\mathbf{q},\mathbf{u}}(\mathbf{X})\otimes \Lambda_{\mathbf{q},\mathbf{u}}(\mathbf{Y})$ denote the ring of symmetric functions in two auxiliary alphabets $\mathbf{X}$ and $\mathbf{Y}$ of infinitely many variables. 
Its homogeneous degree $n$ component has basis $\{s_\lambda(\mathbf{X}) s_\mu(\mathbf{Y})\}$ as $(\lambda, \mu)$ varies over all bipartitions $(\lambda,\mu)$ of $n$. We abbreviate $s_{(\lambda,\mu)} \coloneqq s_\lambda(\mathbf{X}) s_\mu(\mathbf{Y})$.

The \newword{Frobenius character map of type $B$} assigns to each multigraded module 
$M \in \C[\mathfrak{B}_n]$-gmod a symmetric function in 
$\Lambda_{\mathbf{q},\mathbf{u}}(\mathbf{X})\otimes \Lambda_{\mathbf{q},\mathbf{u}}(\mathbf{Y})$, 
such that to each irreducible module $\Frob\left(V^{(\lambda,\mu)}\right) = s_{(\lambda,\mu)}$. 
The Hall inner product is extended by 
\[
    \left\langle f_1(\mathbf{X}) g_1(\mathbf{Y}), f_2(\mathbf{X}) g_2(\mathbf{Y}) \right\rangle 
    = \langle f_1,f_2\rangle \langle g_1, g_2\rangle,
\]
where $\langle f_1,f_2\rangle$ and $\langle g_1, g_2\rangle$ denote 
the usual (type A) Hall inner product on $\Lambda_{\mathbf{q},\mathbf{u}}(\mathbf{X})$ 
and $\Lambda_{\mathbf{q},\mathbf{u}}(\mathbf{Y})$, respectively.
In particular, on this new inner product, the basis $\left\{s_{(\lambda,\mu)}\right\}$ 
is orthonormal, 
with $\left\langle \Frob(M), s_{(\lambda,\mu)} \right\rangle$ counting 
the multiplicity of $V^{(\lambda, \mu)}$ in a given $\mathfrak{B}_n$-module $M$.

The \newword{determinant representation} is $V^{(\varnothing, (1^n))}$ and its associated character is the \newword{determinant character} $\det(\epsilon, \sigma) = \sgn(\sigma) \prod_{i=1}^n \epsilon_i$. This is the type $B$ version of the sign character.
There is an isomorphism $V^{(\lambda,\mu)} \otimes \det \cong V^{(\mu^t,\lambda^t)}$. Every irreducible representation of $\mathfrak{B}_n$ is self-dual.

Let $e_1, \ldots, e_n$ be the standard basis of $\C^n$, on which $\mathfrak{B}_n$ acts by $g \cdot e_i = \epsilon_i e_{\sigma(i)}$, and as before let $e_\tau = e_{\tau_1} \wedge \cdots \wedge e_{\tau_c}$ for any subset $\tau = \{ \tau_1 < \cdots <  \tau_c \} \subseteq [n]$. For $0 \leq c \leq n$, the representation $\bigwedge^c \C^n \cong V^{((n-c),(1^c))}$ is irreducible. 
We call those bipartitions of the form $((n-c),(1^c))$ the \newword{hooks of type $B$}.
For any $\beta \subseteq [n]$ with $|\beta| = c$ we have 
$\bigwedge^c \C^n \cong 
\Ind_{\mathfrak{B}_{\overline{\beta}} \times \mathfrak{B}_\beta}^{\mathfrak{B}_n} 
\left(\mathbf{1}_{\mathfrak{B}_{\overline{\beta}}} \boxtimes \det_{\mathfrak{B}_\beta}\right)$, 
with basis $\{ e_\tau \mid |\tau| = c\}$. 
In particular $\bigwedge^n \C^n = \C \omega$ affords the $1$-dimensional determinant representation, where we write $\omega = e_1 \wedge \cdots \wedge e_n$.

For $\tau, \mu \subseteq [n]$ and $\overline{\mu}=[n]\setminus \mu$, with $|\tau| = |\mu| = c$, define $ \langle e_\tau, e_{\overline{\mu}}\rangle \in \{-1,0,1\}$ by $e_\tau \wedge e_{\overline{\mu}} = \langle e_\tau, e_{\overline{\mu}}\rangle \omega$ 
and extend to a  bilinear pairing $\langle \cdot, \cdot \rangle: \bigwedge^c \C^n \times \bigwedge^{n-c} \C^n \to \bigwedge^n \C^n \cong \C$, so that $\eta \wedge \zeta \mapsto \langle \eta, \zeta \rangle \omega$. 
Since $g \omega = \det(g) \omega$, it satisfies $\langle g \eta, g \zeta\rangle = \det(g) \langle \eta, \zeta\rangle$ for all $g \in \mathfrak{B}_n$, which is the $\mathfrak{B}_n$-equivariant version of Lemma~\ref{lem:pairing-equivariant}.
Then, by the same argument as in Proposition~\ref{prop:Hom-iso}, since
$\left( M \otimes \bigwedge\nolimits^c \C^n \right)^{\det} \cong 
\Hom_{\mathfrak{B}_n} \left(\det, M \otimes \bigwedge\nolimits^c \C^n\right)$, 
for any $\mathfrak{B}_n$-module $M$ and any $0 \leq c \leq n$, there is an isomorphism
\begin{equation}\label{eq:Hom-iso-B}
   \left( M \otimes \bigwedge\nolimits^c \C^n\right)^{\det} 
    \to \Hom_{\mathfrak{B}_n} \left(\bigwedge\nolimits^{n-c} \C^n, M\right), 
    \qquad \text{given by} \qquad f \mapsto \Psi_{f}.
\end{equation}

For $\beta \subseteq [n]$ define type $B$ antisymmetrizers and symmetrizers by
\begin{equation*}
    \cY^B_\beta \coloneqq \frac{1}{|\mathfrak{B}_\beta|} 
    \sum_{g \in \mathfrak{B}_\beta} \det(g) g 
    \qquad \text{and} \qquad 
    \cC^B_{\overline{\beta}} \coloneqq 
    \frac{1}{|\mathfrak{B}_{\overline{\beta}}|} \sum_{g \in \mathfrak{B}_{\overline{\beta}}} g.
\end{equation*}

Then as in Proposition~\ref{prop:hook-hom-iso} for any $\beta \subseteq [n]$ with $|\beta|=n-c$, 
by Frobenius reciprocity, the map $\psi \mapsto \psi(e_\beta)$ induces an isomorphism 
\begin{equation}\label{eq:hoom-hom-iso-B} 
    \Hom_{\mathfrak{B}_n}\left(\bigwedge\nolimits^{n-c} \C^n, M\right)
    \cong \cY^B_\beta \cC^B_{\overline{\beta}}(M)
\end{equation}
for any $\mathfrak{B}_n$-module $M$. 
In particular, observe that in type $B$ there is only 
one irreducible component involved as opposed to two in type $A$. 

\subsection{Bosonic-fermionic coinvariant rings in type B}

Define the diagonal action of each $g \in \mathfrak{B}_n$ on $\C_n^{(k,\ell)}$ by 
\begin{equation*}
    g \cdot x_i^{(\bullet)} = \epsilon_i x_{\sigma(i)}^{(\bullet)} \qquad \text{and} \qquad g \cdot \theta_i^{(\star)} = \epsilon_i \theta_{\sigma(i)}^{(\star)}.
\end{equation*}
The \newword{$(k,\ell)$-bosonic-fermionic coinvariant ring of type $B$} is
\begin{equation}
    \label{eq:coinv-B} 
    R_{\mathfrak{B}_n}^{(k,\ell)} \coloneqq \C_n^{(k,\ell)}\big/ I_{\mathfrak{B}_n}^{(k,\ell)}, \qquad \text{where} \qquad  I_{\mathfrak{B}_n}^{(k,\ell)} \coloneqq 
    \left\langle \left(\C_n^{(k,\ell)}\right)^{\mathfrak{B}_n}_+
    \right\rangle
\end{equation}
is the defining ideal, which is generated by the positive degree polynomials 
invariant under the diagonal action of $\mathfrak{B}_n$. 

Let
$\alpha \subseteq [n]$ and define $\mathfrak{B}_\alpha \subseteq \mathfrak{B}_n$ 
as the subgroup of signed permutations of $\alpha$ which fixes the coordinates in 
$\overline{\alpha} = [n] \setminus \alpha$, so that 
$\mathfrak{B}_\alpha \times \mathfrak{B}_{\overline{\alpha}}$ 
is a subgroup of $\mathfrak{B}_n$. Define $\C_\alpha^{(k,\ell)}$ and 
$p_{\r;\s}(\alpha)$ as in Section~\ref{sec:background}, and set 
$I_{\mathfrak{B}_\alpha}^{(k,\ell)} \coloneqq 
\left\langle \left( \C_\alpha^{(k,\ell)} \right)^{\mathfrak{B}_\alpha}_+ \right\rangle$.

\begin{proposition}[{\cite{JiangLentfer}*{Prop 2.2}}]\label{prop:Weyl-wreath}
Let $n \geq 1$, $k,\ell \geq 0$, and $\alpha \subseteq [n]$.
The (nonunital) algebra $(\C^{(k,\ell)}_\alpha)_+^{\mathfrak{B}_\alpha}$, and thus the ideal
$I_{\mathfrak{B}_\alpha}^{(k,\ell)}$,
is generated by the polarized power sums 
\begin{equation}\label{eq:gen-set-wreath}
    \left \{ p_{\r;\s}(\alpha) \mid
    \r \in (\Z_{\geq 0})^k,\ \s \in \{0,1\}^\ell, 
    2 \leq |\r|+|\s| \leq 2|\alpha|, |\r| + |\s| \text{ even} \right \}.
\end{equation}
In particular, the set~\eqref{eq:gen-set-wreath} at $\alpha = [n]$ 
generates the defining ideal $I_{\mathfrak{B}_n}^{(k,\ell)}$.
\end{proposition}

We note that whenever $|\r| + |\s|$ is even and positive then 
$p_{\r;\s}$ is in $I_{\mathfrak{B}_n}^{(k,\ell)}$.

By Proposition~\ref{prop:Weyl-wreath}, $I_{\mathfrak{B}_n}^{(0,1)} = 0$. 
Thus $R_{\mathfrak{B}_n}^{(0,1)} = \bigwedge[\theta_1, \ldots, \theta_n] \cong \bigwedge \C^n$ via the assignment $\theta_i \mapsto e_i$, with $\left(R_{\mathfrak{B}_n}^{(0,1)}\right)_c = \bigwedge^c \C^n$ for $0 \leq c\leq n$ \cite{JiangLentfer}*{Prop.~3.3}.

\begin{definition} For any $k,\ell \geq 0$, the \newword{multigraded Frobenius series} of $R_{\mathfrak{B}_n}^{(k,\ell)}$ is defined as
    \[
        \Frob\left(R_{\mathfrak{B}_n}^{(k,\ell)}; \mathbf{q}; \mathbf{u}\right) \coloneqq 
        \sum_{(\r, \s) 
        \in (\mathbb{Z}_{\geq 0})^{k} \times \{0,\ldots,n\}^{\ell}}
        \Frob\left(R_{\mathfrak{B}_n}^{(k,\ell)}\right)_{(\r, \s)} 
        q_1^{r_1}\cdots q_k^{r_k} u_1^{s_1}\cdots u_\ell^{s_\ell}, 
    \]
    where $\mathbf{q} = (q_1, \ldots, q_k)$ and $\mathbf{u} = (u_1, \ldots, u_{\ell})$
    are tuples of variables.
\end{definition}

\subsection{Extending the sign isomorphism to type B}\label{subsec:type-B-main}

As in Section~\ref{sec:background}, the map $\tilde{\pi}: \C^{(k,\ell)}_n \otimes \C^{(0,1)}_n \to \C^{(k,\ell+1)}_n$
commutes with the action of $\mathfrak{B}_n$ and induces a multidegree-preserving $\mathfrak{B}_n$-equivariant surjection $\pi:R_{\mathfrak{B}_n}^{(k,\ell)} \otimes R_{\mathfrak{B}_n}^{(0,1)} \to R_{\mathfrak{B}_n}^{(k,\ell+1)}$. We will now show that $\pi$ is in fact an isomorphism on $\det$-isotypic components, thus extending Theorem~\ref{thm:iso-sign} from $\S_n$ to $\mathfrak{B}_n$. 

Before commencing, we note that the proof follows from that of 
Theorem~\ref{thm:iso-sign} mutatis mutandis: use $\mathfrak{B}_n$ (of rank $n$) 
in place of $\mathfrak{S}_n$ (of rank $n-1$), $\det$ in place of $\sgn$, 
and $\bigwedge\C^n$ in place of $\bigwedge\C^{n-1}$.
By comparing Proposition~\ref{prop:Weyl-wreath} with 
Proposition~\ref{prop:ideal-generators}, note that only even-degree 
polarized power sums generate $I_{\mathfrak{B}_n}^{(k,\ell)}$. 
We highlight the main changes to the proof below for the interested reader.

A sequence of signs $\epsilon \in \{\pm 1\}^\beta \subseteq \mathfrak{B}_\beta$ 
acts on $\m(\beta) = m_1(\beta_1) \cdots m_c (\beta_c)$ by $\prod_a \epsilon_{\beta_a}^{\deg m_a}$, 
where $\deg m_{\r;\s} \coloneqq |\r|+|\s|$ is the total degree 
of the monomial with a single index. We have $\det(\epsilon) = \prod_a \epsilon_{\beta_a}$. 
Thus
\begin{equation}\label{eq:YB-monomial}
    \cY^B_\beta ( \m(\beta)) = 
    \begin{cases}
        \cY_\beta ( \m(\beta)) & \text{if all $\deg m_a$ odd,}\\
        0 & \text{otherwise.}
    \end{cases}
\end{equation} 

\begin{definition} 
Let $O_c$ denote the set of all products $\m=m_1\cdots m_c$ where each single index monomial $m_a$ has odd degree for all $1 \leq a \leq c$. For each $\m \in O_c$ define the element $\psi_{\m} \in \Hom_{\mathfrak{B}_n}\left(\bigwedge^c \C^n, R_{\mathfrak{B}_n}^{(k,\ell)}\right)$ whose action on the basis is given by
\[
    \psi_{\m}(e_\tau) \coloneqq \cY_\tau(\m(\tau)).
\]
\end{definition}

It follows from the definition that $\psi_{\m}$ is 
$\mathfrak{B}_n$-equivariant since any sign change affects $e_i$ and $m_a(i)$ in the same way.
We claim the set $\left\{\psi_{\m} \mid \m \in O_{n-c}\right\}$ spans 
$\Hom_{\mathfrak{B}_n}\left(\bigwedge^{n-c} \C^n, R_{\mathfrak{B}_n}^{(k,\ell)}\right)$. 
The proof follows the outline of Theorem~\ref{thm:spanning} 
with some minor changes which we explain below. 

\begin{lemma}\label{lem:spanning-B}
    Suppose $0 \leq c \leq n$ and fix $\beta \subseteq [n]$ with $|\beta| =n-c$. 
    The subspace 
    $\cY^B_\beta\cC^B_{\overline{\beta}}\left(R_{\mathfrak{B}_n}^{(k,\ell)}\right)$
    is spanned by the set
    $\left\{ \psi_{\m}(e_\beta) \mid \m \in O_{n-c}\right\}$. 
    Thus, 
    $\Hom_{\mathfrak{B}_n}\left(\bigwedge^{n-c} \C^n, R_{\mathfrak{B}_n}^{(k,\ell)}\right)$
    is spanned by the set 
    $\left\{ \psi_{\m} \mid \m \in O_{n-c}\right\}$.
\end{lemma}

\begin{proof}
When $c=n$ the subspace consists of constants, so the result holds trivially. So assume $c<n$.
    Consider an element $r \in \cY^B_\beta\cC^B_{\overline{\beta}}\left(R_{\mathfrak{B}_n}^{(k,\ell)}\right)$. 
    Since $I_{\mathfrak{B}_n}^{(k,\ell)}$ is invariant under the action of $\mathfrak{B}_n$, applying $\cY^B_\beta\cC^B_{\overline{\beta}}$ to any representative of $r$ gives a representative $\tilde{r} \in \cY^B_\beta\cC^B_{\overline{\beta}}\left(\C_n^{(k,\ell)}\right)$. 
    As in the proof of Theorem~\ref{thm:spanning}, write $\tilde{r} = \sum_a v_a(\beta)\, u_a(\overline{\beta})$, where $u_a(\overline{\beta})$ is a $\mathfrak{B}_{\overline{\beta}}$-invariant polynomial in the variables indexed by $\overline{\beta}$ which we may assume has no constant term.
    Hence, by Proposition~\ref{prop:Weyl-wreath}, each $u_a(\overline{\beta})$ is a polynomial in the $p_{\r;\s}(\overline{\beta})$ with $|\r| + |\s|$ even and positive, and for these 
    \begin{equation*}
        p_{\r;\s}(\overline{\beta}) = p_{\r;\s}([n]) - p_{\r;\s}({\beta}) \equiv -p_{\r;\s}({\beta}) \bmod I_{\mathfrak{B}_n}^{(k,\ell)},
    \end{equation*}
    because $p_{\r;\s}([n]) \in I_{\mathfrak{B}_n}^{(k,\ell)}$.
    Thus $\tilde{r} \equiv \sum_b c_b\, \m_b(\beta) \bmod{I_{\mathfrak{B}_n}^{(k,\ell)}}$ for some scalars $c_b$ and products $\m_b$ of single-index monomials.
    Applying $\cY^B_\beta$ and using equation~\eqref{eq:YB-monomial}, we write
    \begin{equation*}
        r = \cY^B_\beta(r) = \sum_b c_b\, \cY^B_\beta(\m_b(\beta)) = \sum_{ \substack{ b \text{ such that} \\ \m_b \in O_{n-c}}
        } c_b\, \psi_{\m_b}(e_\beta),
    \end{equation*}
    from which the claims follow. 
\end{proof}

The analogous result to Lemma~\ref{lem:Zero-sum} holds in the following form. Namely, for any $0\leq d \leq n-1$, any $\psi \in \Hom_{\mathfrak{B}_n}\left(\bigwedge^{n-d}\C^n, R_{\mathfrak{B}_n}^{(k,\ell)}\right)$, any $\zeta \in \bigwedge^{n-d-1}\C^n$, and any $(\r,\s)$ with $|\r| + |\s|$ odd, 
    we have $\sum_{i=1}^n \psi(e_i \wedge \zeta) m_{\r;\s}(i) = 0$. 
    This follows the same proof, where $\sum_{i=1}^n m_a(i) m_{\r;\s}(i) = \pm p_{\r^{(a)}+\r; \s^{(a)}+\s}$ is zero or a polarized power sum of even positive degree, so lies in the ideal $I_{\mathfrak{B}_n}^{(k,\ell)}$ again by Proposition~\ref{prop:Weyl-wreath}.

\begin{theorem}\label{thm:iso-det}
    For any $n\geq 1$ and $k, \ell \geq 0$, the  map $\pi: R_{\mathfrak{B}_n}^{(k,\ell)} \otimes R_{\mathfrak{B}_n}^{(0,1)} \to R_{\mathfrak{B}_n}^{(k,\ell+1)}$ is an isomorphism on $\det$-isotypic components. 
\end{theorem}

\begin{proof}
    Let $K_{\mathfrak{B}_n}^{(k,\ell+1)} \coloneqq \left\langle p_{\r;\s;1} : |\r| + |\s| \text{ odd}\right\rangle \subseteq \C_n^{(k,\ell+1)}$. By Proposition~\ref{prop:Weyl-wreath}, the ideal $I_{\mathfrak{B}_n}^{(k,\ell+1)}$ is generated by those $p_{\r;\s;\varepsilon}$ with $|\r| + |\s| + \varepsilon$ even and positive. The polarized power sums with $\varepsilon = 0$ generate $I_{\mathfrak{B}_n}^{(k,\ell)} \otimes \C_n^{(0,1)}$, while those with $\varepsilon = 1$ have $|\r| + |\s|$ odd. 
    Thus Lemma~\ref{lem:Ideal-decomposition} is replaced by $I_{\mathfrak{B}_n}^{(k,\ell+1)} = I_{\mathfrak{B}_n}^{(k,\ell)} \otimes \C_n^{(0,1)} + K_{\mathfrak{B}_n}^{(k,\ell+1)}$.
    The same arguments of Lemma~\ref{lem:tensor-quotient-is-quotient-tensor} and Corollary~\ref{cor:kernel-isotypic} establish that for each irreducible representation of $\mathfrak{B}_n$ indexed by a bipartition $\nu$, 
    the $\nu$-isotypic subspace $\ker(\pi)^\nu$ is spanned by products $f \cdot p_{\r;\s;1}$ with $|\r| + |\s|$ odd and $f \in \left(R_{\mathfrak{B}_n}^{(k,\ell)} \otimes R_{\mathfrak{B}_n}^{(0,1)}\right)^\nu$. In particular, the $R_{\mathfrak{B}_n}^{(0,1)}$ degree $c>0$ component of $\ker(\pi)^{\det}$ is spanned by
    \[
\left\{ f \cdot p_{\r;\s;1} : |\r| + |\s| \equiv 1 \ (\text{mod } 2) \text{ and } f \in \left(R_{\mathfrak{B}_n}^{(k,\ell)} \otimes \bigwedge\nolimits^{c-1}\C^n\right)^{\det}\right\}.
    \]
Since $\bigwedge\nolimits^0 \C^n \cong \C$, then in degree zero the map $\pi$ is an injection, and thus an isomorphism on $\det$-isotypic components. Thus, we may assume $c>0$. 
    
Hence, by equation \eqref{eq:Hom-iso-B} in order to show $f \cdot p_{\r;\s;1} =0$ it suffices to prove that $\Psi_{f \cdot p_{\r;\s;1}} =0$ for all $ f \in \left(R_{\mathfrak{B}_n}^{(k,\ell)} \otimes \bigwedge\nolimits^{c-1}\C^n\right)^{\det}$, with $\Psi_{f \cdot p_{\r;\s;1}} \in \Hom_{\mathfrak{B}_n} \left(\bigwedge\nolimits^{n-c} \C^n, R_{\mathfrak{B}_n}^{(k,\ell)}\right)$.

Indeed, for any such $f$ and $\mu \subseteq [n]$ with $|\mu| = c$, the same proof as in 
Theorem~\ref{thm:iso-sign} yields the identity 
   \begin{equation*}
        \Psi_{f \cdot p_{\r;\s;1}}(e_{\overline{\mu}}) = \pm \sum_{i=1}^n \Psi_f(e_i \wedge e_{\overline{\mu}})\, m_{\r;\s}(i),
    \end{equation*}
    with a sign depending only on $c$ and $|\s|$ and $\Psi_f \in 
    \Hom_{\mathfrak{B}_n} \left(\bigwedge\nolimits^{n+1-c} \C^n, 
    R_{\mathfrak{B}_n}^{(k,\ell)}\right)$. 
    Then, by the analogue of Lemma \ref{lem:Zero-sum} discussed above 
    evaluated at $d=c-1$, the sum 
    $\sum_{i=1}^n \Psi_f(e_i \wedge e_{\overline{\mu}})\, m_{\r;\s}(i)$, 
    and thus $\Psi_{f \cdot p_{\r;\s;1}}(e_{\overline{\mu}})$ 
    are zero for all $\mu$. Therefore, $\ker(\pi)^{\det}$ is trivial 
    and thus $\pi$ is an isomorphism on $\det$-isotypic components. 
\end{proof}

\begin{corollary}\label{cor:hook-to-sign-sf-B}
    For any $n \geq 1$ and $k, \ell \geq 0$ we have
    \begin{equation*}
    \begin{aligned}
        \left\langle 
        \Frob\left(R_{\mathfrak{B}_n}^{(k,\ell+1)}; 
        \mathbf{q}; \mathbf{u}, a\right), s_{(\varnothing, (1^n))} 
        \right\rangle
        &= \left\langle 
        \Frob\left(R_{\mathfrak{B}_n}^{(k,\ell)} \otimes R_{\mathfrak{B}_n}^{(0,1)}; 
        \mathbf{q}; \mathbf{u}, a\right), s_{(\varnothing,(1^n))} 
        \right\rangle \\
        &= \sum_{d=0}^{n} a^d \left\langle 
        \Frob\left(R_{\mathfrak{B}_n}^{(k,\ell)}; \mathbf{q}; \mathbf{u}\right), 
        s_{((d),(1^{n-d}))} \right\rangle. 
    \end{aligned}
    \end{equation*}
\end{corollary}

\begin{proof}
    The first equality follows from Theorem~\ref{thm:iso-det} by applying the type $B$ Frobenius character map to the $\det$-isotypic components, multidegree by multidegree.

    By self-duality of $\mathfrak{B}_n$-modules, and since $\det \otimes \bigwedge^c \C^n \cong \bigwedge^{n-c} \C^n$, both sides of the isomorphism of equation~\eqref{eq:Hom-iso-B} are the multiplicity space of $V^{((c),(1^{n-c}))}$ in a $\mathfrak{B}_n$-module $M$. 
    The multiplicity of $\det$ in $R_{\mathfrak{B}_n}^{(k,\ell)} \otimes \bigwedge^c \C^n$ is equal to that of $V^{((c),(1^{n-c}))}$ in $R_{\mathfrak{B}_n}^{(k,\ell)}$.
    Taking the Frobenius image and summing over the multidegrees completes the proof of the second equality.
\end{proof}

\section{Combinatorial consequences}\label{sec:consequences}

By Corollary~\ref{cor:hook-to-sign-sf}, we can determine the multiplicity of the sign character of $R_n^{(k,\ell+1)}$ as long as we know all of the hook characters of $R_n^{(k,\ell)}$. These hook characters are known for all $k+\ell \leq 2$: for $R_n^{(1,0)}$ see \cite{stanley-invariants-finite-groups}, for $R_n^{(0,1)}$ see \cite{HaglundSergel}, for $R_n^{(2,0)}$ see \cites{Haglund2008,Haglund2004,Alfano,Haiman2002}, for $R_n^{(1,1)}$ see \cite{MuraiRhoadesWilson}, and for $R_n^{(0,2)}$ see \cite{KimRhoades2022}.  
This implies that the existing results on the sign characters of $R_n^{(1,1)}$ \cite{SwansonWallach1}, $R_n^{(0,2)}$ \cite{KimRhoades2022}, $R_n^{(2,1)}$ \cite{GLM}, and $R_n^{(0,3)}$ \cite{Lentfer2025} may be recovered.\footnote{For $R_n^{(0,3)}$, there is an erroneous extra term of $uv[n-1]_{u,v}$ in the published version of \cite{Lentfer2025}*{Prop. 9}, which should say that $\langle \Frob(R_n^{(0,2)} \otimes R_n^{(0,1)};u,v,w),s_{(1^n)} \rangle = s_{(n-1)}(u,v,w) + s_{(n-2,1,1)}(u,v,w)$. Then this agrees with Corollary~\ref{cor:hook-to-sign-sf}.}  
The sign character of $R_n^{(1,2)}$ had not previously been established, which we do shortly.

Similarly, in type $B$, by Corollary~\ref{cor:hook-to-sign-sf-B}, we can determine the multiplicity of the determinant character of $R_{\mathfrak{B}_n}^{(k,\ell+1)}$ as long as we know all of the hook characters of $R_{\mathfrak{B}_n}^{(k,\ell)}$.
These hook characters are known for several small cases: for $R_{\mathfrak{B}_n}^{(1,0)}$ see \cite{Stembridge}, for $R_{\mathfrak{B}_n}^{(0,1)}$ see \cite{JiangLentfer}, for $R_{\mathfrak{B}_n}^{(1,1)}$ see \cites{BR, BCJLR}, and for $R_{\mathfrak{B}_n}^{(0,2)}$ see \cites{KimRhoades2022, JiangLentfer}. 
This implies that the existing results on the determinant characters of $R_{\mathfrak{B}_n}^{(1,1)}$ \cite{SwansonWallach1}, $R_{\mathfrak{B}_n}^{(0,2)}$ \cite{KimRhoades2022}, and $R_{\mathfrak{B}_n}^{(0,3)}$ \cite{JiangLentfer} may be recovered. 
The determinant character of $R_{\mathfrak{B}_n}^{(1,2)}$ had not previously been established, which we do shortly.
In contrast with type $A$, the hook characters are not known for $R_{\mathfrak{B}_n}^{(2,0)}$.

\subsection{Sign and determinant characters of \texorpdfstring{$(1,2)$}{(1,2)}-bosonic-fermionic coinvariant rings}\label{subsec:sign-1-2}

\begin{proposition}[{\cite{BCJLR}*{Prop. 4.10}}]\label{prop:bcjlr-A}
    For $n \geq 1$,
    \begin{equation*}
        \left\langle 
        \Frob \left(R_n^{(1,1)} \otimes R_n^{(0,1)}; q; u, v\right), s_{(1^n)} 
        \right\rangle = 
        \sum_{0 \leq i+j \leq n-1} u^i v^j q^{\binom{n-i-j}{2}} \qbinom{n-1}{j}_q\qbinom{n-1-j}{i}_q.
    \end{equation*}
\end{proposition}

We deduce the following, whose ungraded version was conjectured by Bergeron \cite{Bergeron2020}.

\begin{corollary}\label{cor:12-sign}
For $n \geq 1$,
    \begin{equation*}
        \left\langle 
        \Frob \left(R_n^{(1,2)}; q; u, v\right), s_{(1^n)} \right\rangle = 
        \sum_{0 \leq i+j \leq n-1} u^i v^j q^{\binom{n-i-j}{2}} \qbinom{n-1}{j}_q\qbinom{n-1-j}{i}_q.
    \end{equation*}
Consequently, the multiplicity of the sign character of $R_n^{(1,2)}$ is $3^{n-1}$.
\end{corollary}

\begin{proof}
    Immediate from combining Corollary~\ref{cor:hook-to-sign-sf} with Proposition~\ref{prop:bcjlr-A}.
\end{proof}

\begin{proposition}[{\cite{BCJLR}*{Prop. 4.5}}]\label{prop:bcjlr-B}
    For $n \geq 1$,
    \begin{equation*}
        \left\langle 
        \Frob \left(R_{\mathfrak{B}_n}^{(1,1)} \otimes R_{\mathfrak{B}_n}^{(0,1)}; q; u, v\right), 
        s_{(\varnothing,(1^n))} \right\rangle = 
        \sum_{d=0}^n v^d \qbinom{n}{d}_{q^2} \prod_{j=1}^{n-d} (u+q^{2j-1}). 
    \end{equation*}
\end{proposition}

We deduce the following, whose ungraded version was conjectured by Jiang and the second author \cite{JiangLentfer}.
The trigraded version was shown to be an upper bound in \cite{BCJLR}, and now we upgrade it to an equality.
\begin{corollary}\label{cor:12-sign-B}
For $n \geq 1$,
    \begin{equation*}
        \left\langle 
        \Frob \left(R_{\mathfrak{B}_n}^{(1,2)}; q; u, v\right), 
        s_{(\varnothing,(1^n))} \right\rangle = 
        \sum_{d=0}^n v^d \qbinom{n}{d}_{q^2} \prod_{j=1}^{n-d} (u+q^{2j-1}). 
    \end{equation*}
Consequently, the multiplicity of the determinant character of $R_{\mathfrak{B}_n}^{(1,2)}$ is $3^{n}$.
\end{corollary}

\begin{proof}
    Immediate from combining Corollary~\ref{cor:hook-to-sign-sf-B} with Proposition~\ref{prop:bcjlr-B}.
\end{proof}

%%%%%%%%%%%%%%%%%%%%%%%%%%%%%%%%%%%%%%%%%%%%%%%%%%%%%%%%%%%%%%%%%%%%%%
\subsection{A proof of the hook components conjecture of Bergeron}
%%%%%%%%%%%%%%%%%%%%%%%%%%%%%%%%%%%%%%%%%%%%%%%%%%%%%%%%%%%%%%%%%%%%%%

In this section, we deduce a conjecture of Bergeron \cite{BergeronOPAC}.

Define a \newword{super Schur function} (see \cites{BereleRegev,ChengWangBook}) by 
\begin{equation}\label{eq:superschur}
    s_\lambda(\mathbf{q}/\mathbf{u}) = \sum_{\nu \subseteq \lambda} s_\nu(\mathbf{q}) s_{\lambda^t/\nu^t}(\mathbf{u}).
\end{equation}
We note that super Schur functions may also be defined in terms of certain super tableaux (see \cite{Musson}*{Chap. 12.3}). 
Define a \newword{vertical strip} to be a skew shape $\lambda/\rho$ with no two boxes in the same row (see \cite{StanleyEC2}*{Chap. 7.15}). 
Define $P(k,\ell,n)$ to be the set of all partitions with length $l (\lambda) \leq n$ and $\lambda_{k+1} \leq \ell$.
We recall diagonal supersymmetry, conjectured in \cite{Bergeron2020}*{Conj. 1}.
\begin{theorem}[{\cite{Lentfer-Supersymmetry}*{Cor. 1.3}}]\label{thm:DSS-S_n}
    Fix a positive integer $n$. 
    For partitions $\lambda$ with $l(\lambda) \leq n$, and $\mu \vdash n$, there exist nonnegative integer coefficients $c_{\lambda,\mu}$ such that for any $(k,\ell)$, the multigraded Frobenius series of $R_n^{(k,\ell)}$ is
\begin{equation}\label{eq:coefficients-def}
    \Frob\left(R_n^{(k,\ell)}; \mathbf{q};\mathbf{u}\right) = 
    \sum_{\lambda \in P(k,\ell,n)} \sum_{\mu \vdash n} 
    c_{\lambda,\mu} s_\lambda(\mathbf{q}/\mathbf{u})s_\mu. 
\end{equation}
\end{theorem}

We now prove the hook components conjecture of Bergeron \cite{BergeronOPAC}*{Conj. 2}, 
which was formulated with the goal of reconstructing the Hilbert series of 
$\left(R_n^{(k,\ell)}\right)^{\sgn}$. 
Recall that the universal coefficients $c_{\lambda,\mu}$ 
are defined by equation~\eqref{eq:coefficients-def}. 
Furthermore, define $c_{\lambda,\mu}\coloneqq0$ if $l(\lambda)>n$. 

\begin{corollary}[\cite{BergeronOPAC}*{Conj. 2 (Hook components)}]\label{cor:hook-components} 
Fix $n\geq1$. 
For all $0\leq d\leq n-1$ and all partitions $\rho$,
\begin{equation}\label{eq:hook-components}
    c_{\rho,(d+1,1^{n-d-1})} = \sum_{\lambda} c_{\lambda,(1^n)},
\end{equation}
where the sum is over all partitions $\lambda\supseteq\rho$ such that $\lambda/\rho$ is a vertical strip of size $d$.
\end{corollary}

\begin{proof}
    Fix $K\geq n$.
    By Theorem~\ref{thm:DSS-S_n}, and the definition of super Schur functions, we have that
    \begin{equation*}
            \left\langle 
            \Frob \left(R_n^{(K,1)}; \mathbf{q}; a\right), s_{(1^n)} 
            \right\rangle = 
            \sum_{\lambda} c_{\lambda,(1^n)}\, s_\lambda(\mathbf{q}/a) = 
            \sum_{d\geq0} a^d \sum_{\rho}\left(\sum_{\lambda} c_{\lambda,(1^n)}\right) 
            s_\rho(\mathbf{q}), 
    \end{equation*}
    where the inner sum is over the partitions $\lambda\supseteq\rho$ such that $\lambda/\rho$ is a vertical strip of size $d$.

    We also have that $\left\langle \Frob \left(R_n^{(K,0)}; \mathbf{q}\right), 
    s_{(d+1,1^{n-d-1})} \right\rangle = 
    \sum_\rho c_{\rho,(d+1,1^{n-d-1})}\,s_\rho(\mathbf{q})$. 
    The Schur polynomials $s_\rho(q_1,\ldots,q_K)$ with $l(\rho)\leq n\leq K$ are linearly independent \cite{Macdonald}, so comparing the coefficients of $a^d s_\rho(\mathbf{q})$ on the two sides of Corollary~\ref{cor:hook-to-sign-sf} at $(K,0)$ proves the claim when $l(\rho)\leq n$. 
    If $l(\rho)>n$, both sides are zero (with the extended definition of the $c_{\lambda,\mu}$).
\end{proof}

\begin{remark}\label{rem:bergeron-modules}
In \cite{BergeronOPAC}, the module $\mathcal{E}_n^{\langle k \rangle}$ is introduced as the module of $k$-variate diagonal harmonic polynomials (building on \cite{Bergeron2013}), described as the smallest $\S_n$-submodule of $\C_n^{(k,0)}$ which contains the Vandermonde determinant in the variables $\bm{x}^{(1)}$ and is closed under partial derivatives and under higher polarization operators 
\begin{equation*}
    \sum_{i=1}^n x_i^{(a)} \partial_{x_i^{(b)}}^{m} \text{ for $1 \leq a,b \leq k$ and $m \geq 1$}.
\end{equation*}
The module described in this way is contained in the space of diagonal harmonics $H_n^{(k,0)}$, which is isomorphic to $R_n^{(k,0)}$ as a multigraded $\S_n$-module (see \cite{Haiman1994} or \cite{JiangLentfer}*{Appendix A}).
It is easy to see that this construction agrees with the definition of $H_n^{(k,0)}$ in the case of $k=1$, and when $k=2$ they agree by Haiman's operator theorem \cite{Haiman2002}*{Thm. 4.2}.
For $k \geq 3$, they are expected to agree (see \cite{Haiman1994}*{Sec. 5.1}, \cite{Haiman2002}*{Sec. 1}, \cite{BergeronPrevilleRatelle}*{Sec. 2}, and \cite{Zabrocki2020}), but to our knowledge this has not been proven.
In this paper, when we say that we prove Bergeron's hook-components conjecture (in Corollary~\ref{cor:hook-components}), we prove it interpreting $\mathcal{E}_n^{\langle k \rangle}$ to be the entire harmonic space $H_n^{(k,0)}$, which is isomorphic to the coinvariant space $R_n^{(k,0)}$.
\end{remark}

An analogous result holds in type $B$.
By \cite{Lentfer-Supersymmetry}*{Thm.~1.2}, there exist nonnegative integers $c_{\lambda,(\mu^{(1)},\mu^{(2)})}$, for partitions $\lambda$ with $l(\lambda) \leq n$ and bipartitions $(\mu^{(1)},\mu^{(2)})$ of $n$, such that for all $(k,\ell)$,
\begin{equation}\label{eq:coefficients-def-B}
    \Frob\left(R_{\mathfrak{B}_n}^{(k,\ell)}; \mathbf{q}; \mathbf{u}\right) = 
    \sum_{\lambda \in P(k,\ell,n)} \sum_{(\mu^{(1)},\mu^{(2)}) \vdash n} 
    c_{\lambda,(\mu^{(1)},\mu^{(2)})} s_\lambda(\mathbf{q}/\mathbf{u}) s_{(\mu^{(1)},\mu^{(2)})}. 
\end{equation}
Further set $c_{\lambda,(\mu^{(1)},\mu^{(2)})} \coloneqq 0$ if $l(\lambda) > n$.
Now we record the type $B$ version of Corollary~\ref{cor:hook-components}.

\begin{corollary}\label{cor:hook-components-B}
    Fix $n \geq 1$. 
    For all $0 \leq d \leq n$ and all partitions $\rho$,
    \begin{equation}\label{eq:hook-components-B}
        c_{\rho,((d),(1^{n-d}))} = \sum_{\lambda} c_{\lambda,(\varnothing,(1^n))},
    \end{equation}
    where the sum is over all partitions $\lambda \supseteq \rho$ such that $\lambda/\rho$ is a vertical strip of size $d$.
\end{corollary}
 
\begin{proof}
    Fix $K\geq n$.
    By equation~\eqref{eq:coefficients-def-B}, and the definition of super Schur functions, we have that
    \begin{equation*}
            \left\langle 
            \Frob \left(R_{\mathfrak{B}_n}^{(K,1)}; \mathbf{q}; a\right), 
            s_{(\varnothing,(1^n))} \right\rangle = 
            \sum_{\lambda} c_{\lambda,(\varnothing,(1^n))}\, s_\lambda(\mathbf{q}/a) = 
            \sum_{d\geq0} a^d \sum_{\rho}\left(\sum_{\lambda} c_{\lambda,(\varnothing,(1^n))}\right) 
            s_\rho(\mathbf{q}), 
    \end{equation*}
    where the inner sum is over the partitions $\lambda\supseteq\rho$ such that $\lambda/\rho$ is a vertical strip of size $d$.

    We also have that 
    $\left\langle \Frob \left(R_{\mathfrak{B}_n}^{(K,0)}; \mathbf{q}\right), 
    s_{((d),(1^{n-d}))} \right\rangle = 
    \sum_\rho c_{\rho,((d),(1^{n-d}))}\,s_\rho(\mathbf{q})$. 
    The Schur polynomials $s_\rho(q_1,\ldots,q_K)$ with $l(\rho)\leq n\leq K$ are linearly independent, so comparing the coefficients of $a^d s_\rho(\mathbf{q})$ on the two sides of Corollary~\ref{cor:hook-to-sign-sf-B} at $(K,0)$ proves the claim when $l(\rho)\leq n$, while if $l(\rho)>n$, both sides are zero.
\end{proof}

%%%%%%%%%%%%%%%%%%%%%%%%%%%%%%%%%%%%%%%%%%%%%%%%%%%%%%%%%%%%%%%%%%%%%%
\section*{Computational tools and AI usage}

Computational evidence for some of the results was collected using SageMath, Macaulay2, and Python. 
We utilized modifications of the Jupyter notebooks \cite{Zabrocki-code} and \cite{Bergeron-code}. 
Generative AI tools were used to reinterpret constructions in \cite{GLM} 
from harmonic spaces into coinvariant spaces. 
The proofs and text in this paper were entirely written by the authors 
and they take complete responsibility for its contents. 
%%%%%%%%%%%%%%%%%%%%%%%%%%%%%%%%%%%%%%%%%%%%%%%%%%%%%%%%%%%%%%%%%%%%%%

\bibliography{biblio}
\end{document}
